\subsection{Measures defined by numerical densities}

DEFINITION 2. --- A function f defined on T, with values in \(\overline{R}\) or in a Banach space, is said to be locally \(\mu\) -integrable if f is \(\mu\) -measurable and every point \(x \in T\) admits a neighborhood V such that \(\mu^{\bullet}(|\mathbf{f}|\varphi_{\mathrm{V}}) < +\infty\) .

This definition coincides with the one given in Ch. V, §5, No.~1, in case T is locally compact.

Let \(f\) be a locally \(\mu\) -integrable positive function; the mapping \(K \mapsto f_K \cdot \mu_K\) is a premeasure (Ch. V, §7, No.~1, Cor. 2 of Th. 1), which will be denoted \(f \cdot \mu\) .

PROPOSITION 2. --- If \(f\) is a positive locally \(\mu\) -integrable function, then, for every function \(g \in \mathcal{F}_{+}(\mathrm{T})\) , one has the relation

\[
(f \cdot \mu) ^ {\bullet} (g) = \mu^ {\bullet} (f g).\tag{1}
\]

Indeed, for every compact set K in T,

\[
(f \cdot \mu) _ {\mathrm{K}} ^ {\bullet} (g _ {\mathrm{K}}) = (f _ {\mathrm{K}} \cdot \mu_ {\mathrm{K}}) ^ {\bullet} (g _ {\mathrm{K}}) = \mu_ {\mathrm{K}} ^ {\bullet} (f _ {\mathrm{K}} g _ {\mathrm{K}}) = \mu_ {\mathrm{K}} ^ {\bullet} ((f g) _ {\mathrm{K}}),
\]

on using the definition of \(f \cdot \mu\) and Prop. 3 of Ch. V, §5, No.~3. Prop. 2 follows from this on passing to the supremum over \(K\) .

Now let \(x \in \mathbf{T}\) and let \(\mathbf{V}\) be a neighborhood of \(x\) such that \(\mu^{\bullet}(f\varphi_{\mathrm{V}}) < +\infty\) (Def. 2); then \((f \cdot \mu)^{\bullet}(\mathrm{V}) = \mu^{\bullet}(f\varphi_{\mathrm{V}}) < +\infty\) , therefore \(f \cdot \mu\) is a measure.

DEFINITION 3. --- Let \(f\) be a locally \(\mu\) -integrable positive function. The measure \(f \cdot \mu : K \mapsto f_K \cdot \mu_K\) is called the measure with density \(f\) with respect to \(\mu\) , or the product measure of \(\mu\) by the function f.~Every measure of the form \(f \cdot \mu\) , where f is positive and locally \(\mu\) -integrable, is called a measure with base \(\mu\) .

Remarks. --- 1) The definition of \(f \cdot \mu\) extends to the case that \(f\) is a complex locally integrable function; one then has \(|f \cdot \mu| = |f| \cdot \mu\) , which implies at once that \(f \cdot \mu\) is a measure, not just a premeasure. We retain the expression 'measures with base \(\mu'\) to designate the complex measures so defined.

\begin{enumerate}
\def\labelenumi{\arabic{enumi})}
\setcounter{enumi}{1}
\tightlist
\item
  Similarly, if \(\theta\) is a complex measure, \(f\) is said to be locally \(\theta\) -integrable if it is locally \(|\theta|\) -integrable, and one defines the measure \(f \cdot \theta : K \mapsto f_K \cdot \theta_K\) . One has \(|f \cdot \theta| = |f| \cdot |\theta|\) (Ch. V, §5, No.~2, Prop. 2). In this No., we shall leave aside everything concerning non-positive measures.
\end{enumerate}

PROPOSITION 3. --- Let \(\nu\) be a measure on T. For \(\nu\) to be of the form \(f \cdot \mu\) , where \(f\) is a locally \(\mu\) -integrable positive function, it is necessary and sufficient that every \(\mu\) -negligible compact set be \(\nu\) -negligible. If \(f'\) is a second locally \(\mu\) -integrable function such that \(\nu = f' \cdot \mu\) , then \(f = f'\) locally \(\mu\) -almost everywhere.

The condition is obviously necessary (Prop. 2). Conversely, suppose that every \(\mu\) -negligible compact set is \(\nu\) -negligible. Let us introduce a crushing \((\mathrm{K}_{\alpha})_{\alpha \in \mathrm{A}}\) of \(\mathbf{T}\) for the measure \(\mu + \nu\) and let us set \(\mathbf{N} = \mathbf{T} - \bigcup_{\alpha \in \mathrm{A}} \mathrm{K}_{\alpha}\) . It is clear that \((\mathrm{K}_{\alpha})_{\alpha \in \mathrm{A}}\) is a crushing for \(\mu\) and for \(\nu\) , and Prop. 9 of §1, No.~8 therefore implies the following relations for every \(g \in \mathcal{F}_+\) :

\[
\mu^ {\bullet} (g) = \sum_ {\alpha \in A} \mu_ {K _ {\alpha}} ^ {\bullet} \left(g _ {K _ {\alpha}}\right), \quad \nu^ {\bullet} (g) = \sum_ {\alpha \in A} \nu_ {K _ {\alpha}} ^ {\bullet} \left(g _ {K _ {\alpha}}\right).
\]

Consider a compact set \(C \subset K_{\alpha}\) that is \(\mu_{K_{\alpha}}\) -negligible; then C is locally \(\mu\) -negligible, hence locally \(\nu\) -negligible, and finally \(\nu_{K_{\alpha}}\) -negligible by the definition of \(\nu\) . It then follows from the Lebesgue--Nikodym theorem (Ch. V, §5, No.~5, Th. 2) that \(\nu_{K_{\alpha}}\) admits a density \(f_{\alpha}\) with respect to \(\mu_{K_{\alpha}}\) . Let f be the function that coincides with \(f_{\alpha}\) on each of the sets \(K_{\alpha}\) , and with 0 on N; the function f is \(\mu\) -measurable (Ch. IV, §5, No.~10, Prop. 16), and for every function \(g \in F_{+}\) one has, by the above relations and Prop. 3 of Ch. V, §5, No.~3,

\[
\nu^ {\bullet} (g) = \sum_ {\alpha \in \mathrm{A}} \nu_ {\mathrm{K} _ {\alpha}} ^ {\bullet} \left(g _ {\mathrm{K} _ {\alpha}}\right) = \sum_ {\alpha \in \mathrm{A}} \mu_ {\mathrm{K} _ {\alpha}} ^ {\bullet} \left(f _ {\alpha} g _ {\mathrm{K} _ {\alpha}}\right) = \sum_ {\alpha \in \mathrm{A}} \mu_ {\mathrm{K} _ {\alpha}} ^ {\bullet} ((f g) _ {\mathrm{K} _ {\alpha}}) = \mu^ {\bullet} (f g).
\]

It then follows first of all that \(f\) is locally \(\mu\) -integrable: if \(x\) is a point of \(T\) , and if \(V\) is a neighborhood of \(x\) such that \(\nu^{\bullet}(V) < +\infty\) , then \(\mu^{\bullet}(f\varphi_V) < +\infty\) . Next, Prop. 2 shows that the measures \(\nu\) and \(f \cdot \mu\) have the same essential upper integral. They are therefore equal (§1, No.~2, Cor. of Prop. 2). The uniqueness of \(f\) being obvious based on the case of compact spaces, the proposition is established.

Remark 3). --- The theory of integration with respect to a measure \(\nu = f \cdot \mu\) reduces at once to the theory treated in Ch. V. For, let \((\mathrm{K}_{\alpha})_{\alpha \in \mathbf{A}}\) be a crushing of T for \(\mu\) , hence for \(\nu\) , and let \(\mathrm{T}'\) be the locally compact space defined in the Scholium of §1, No.~8; we can associate to \(\mu\) (resp. to \(\nu\) ) a measure \(\mu'\) (resp. \(\nu'\) ) on \(\mathrm{T}'\) , in such a way that the measurable functions, the essentially integrable functions with values in a Banach space, and the essential upper integrals of positive functions, are the same for \(\mu\) and \(\mu'\) (resp. for \(\nu\) and \(\nu'\) ). The function \(f\) is therefore \(\mu'\) -measurable; it is locally \(\mu'\) -integrable, because \(\mathrm{T}'\) is locally compact, and a compact subset of \(\mathrm{T}'\) intersects only a finite number of the compact sets \(\mathrm{K}_{\alpha}\) ( \(\alpha \in \mathbf{A}\) ). Finally, the relation \(\nu'^{\bullet}(g) = \nu^{\bullet}(g) = \mu^{\bullet}(fg) = \mu'^{\bullet}(fg)\) proves that \(\nu' = f \cdot \mu'\) (Ch. V, §5, No.~3, Prop. 3). We leave to the reader the task of transcribing the results of Ch. V, §5.

\subsection{Image of a measure}

DEFINITION 4. --- Let \(\pi\) be a mapping of T into a topological space X. One says that \(\pi\) is \(\mu\) -proper if \(\pi\) is \(\mu\) -measurable, and if every point \(x\) of X admits a neighborhood V such that \(\mu^{\bullet}\left(\overline{\pi}^{-1}(\mathrm{V})\right) < +\infty\) .

Remarks. -- 1) When T and X are locally compact, this definition is equivalent to that of Ch. V, §6, No.~1.

\begin{enumerate}
\def\labelenumi{\arabic{enumi})}
\setcounter{enumi}{1}
\item
  A proper continuous mapping (GT, I, §10, No.~1, Def. 1) of T into X is \(\mu\) -proper for every measure \(\mu\) . For, let \(x \in X\) ; since \(\overline{\pi}^{-1}(x)\) is compact (loc. cit., No.~2, Th. 1), the set \(\overline{\pi}^{-1}(x)\) has an open neighborhood H such that \(\mu^{\bullet}(\mathrm{H}) < +\infty\) . Set \(\mathrm{V} = \mathrm{X} - \pi(\mathrm{T} - \mathrm{H})\) ; since \(\pi\) is closed, V is open in X, contains \(x\) , and satisfies \(\overline{\pi}^{-1}(\mathrm{V}) \subset \mathrm{H}\) , whence \(\mu^{\bullet}(\overline{\pi}^{-1}(\mathrm{V})) \leqslant \mu^{\bullet}(\mathrm{H}) < +\infty\) .
\item
  If \(\mu\) is bounded, every \(\mu\) -measurable mapping of T into X is \(\mu\) -proper.
\item
  If \(\theta\) is a complex measure on \(\mathbf{T}\) , \(\pi\) is said to be \(\theta\) -proper if \(\pi\) is proper for the positive measure \(|\theta|\) .
\end{enumerate}

PROPOSITION 4. --- Let \(\pi\) be a \(\mu\) -proper mapping of T into a topological space X. There exists one and only one measure \(\nu\) on X such that \(\nu^{\bullet}\) is equal to the image encumbrance \(\pi(\mu^{\bullet})\) (§1, No.~1), in other words, such that \(\nu^{\bullet}(g) = \mu^{\bullet}(g \circ \pi)\) for all \(g \in \mathcal{F}_{+}(\mathbf{X})\) .

Uniqueness is obvious (§1, No.~2, Cor. of Prop. 2). To establish existence, we shall first treat the case that \(\mu\) is carried by a compact set K, such that the restriction of \(\pi\) to K is continuous. Then \(L = \pi(K)\) is compact; let \(\pi'\) be the continuous mapping of K into L induced by \(\pi\) , and let \(\nu'\) be the image measure \(\pi'(\mu_K)\) on L, \(\nu\) the measure on X defined by \(\nu'\) (§1, No.~3, Example 2). For every \(g \in \mathcal{F}_+(X)\) ,

\[
\nu^ {\bullet} (g) = \nu^ {\prime \bullet} (g _ {\mathrm{L}}) = \mu_ {\mathrm{K}} ^ {\bullet} (g _ {\mathrm{L}} \circ \pi^ {\prime}) = \mu_ {\mathrm{K}} ^ {\bullet} ((g \circ \pi) _ {\mathrm{K}}) = \mu^ {\bullet} ((g \circ \pi) _ {\mathrm{K}} ^ {0}) = \mu^ {\bullet} (g \circ \pi)
\]

(we have used successively formula (3) of §1, No.~3; Prop. 2 of Ch. V, §6, No.~2; the definition of \(\mu_{\mathrm{K}}^{\bullet}\) ; and the fact that \(\mu\) is carried by K). In other words, \(\nu^{\bullet} = \pi(\mu^{\bullet})\) .

Let us now pass to the general case; by Props. 10 and 9 of §1, No.~8, \(\mu\) is the sum of a summable family \((\mu_{\alpha})_{\alpha \in \mathbf{A}}\) of measures with compact support, such that the restriction of \(\pi\) to the support \(\mathbf{K}_{\alpha}\) of \(\mu_{\alpha}\) is continuous for every \(\alpha \in \mathbf{A}\) . The special case treated above permits associating to each measure \(\mu_{\alpha}\) on \(\mathbf{T}\) a measure \(\nu_{\alpha}\) on \(\mathbf{X}\) such that \(\nu_{\alpha}^{\bullet} = \pi (\mu_{\alpha}^{\bullet})\) . Then, for \(g \in \mathcal{F}_{+}(\mathbf{X})\) ,

\[
\sum_ {\alpha \in A} \nu_ {\alpha} ^ {\bullet} (g) = \sum_ {\alpha \in A} \mu_ {\alpha} ^ {\bullet} (g \circ \pi) = \mu^ {\bullet} (g \circ \pi).
\]

The encumbrance \(\pi (\mu^{\bullet})\) is locally bounded, since \(\pi\) is \(\mu\) -proper; the family \((\nu_{\alpha})_{\alpha \in \mathbf{A}}\) is therefore summable (§1, No.~7, Prop. 7), and its sum \(\nu\) satisfies the statement.

DEFINITION 5. --- If \(\pi\) is a \(\mu\) -proper mapping of T into a topological space X, the unique measure \(\nu\) on X such that \(\nu^{\bullet}(g) = \mu^{\bullet}(g \circ \pi)\) for all \(g \in \mathcal{F}_{+}(\mathrm{X})\) is called the image measure of \(\mu\) under \(\pi\) , and is denoted \(\pi(\mu)\) .

Example. --- Let K be a compact subspace of T, i the canonical injection of K into T, and \(\lambda\) a measure on K; since i is continuous and \(\lambda\) is bounded, i is \(\lambda\) -proper. Formula (3) of §1, No.~3 shows that the ``measure on T defined by \(\lambda\) '' is the image measure \(i(\lambda)\) .

Remark 5). --- If \(\theta\) is a real measure and \(\pi\) is \(\theta\) -proper, then \(\pi\) is proper for the measures \(\theta^{+}\) and \(\theta^{-}\) ; one then sets \(\pi(\theta) = \pi(\theta^{+}) - \pi(\theta^{-})\) . If \(\theta\) is a complex measure and \(\pi\) is \(\theta\) -proper, then \(\pi\) is proper for the real measures \(\mathcal{R}(\theta)\) and \(\mathcal{I}(\theta)\) ; one then sets

\[
\pi (\theta) = \pi (\mathcal {R} (\theta)) + i \pi (\mathcal {I} (\theta)).
\]

PROPOSITION 5. --- Let \(\pi\) be a \(\mu\) -proper mapping of T into a topological space X, and let \(f\) be a mapping of X into a topological space F (Hausdorff or not). For \(f\) to be \(\pi(\mu)\) -measurable, it is necessary and sufficient that \(f \circ \pi\) be \(\mu\) -measurable.

Let us take up again the proof of Prop. 4, and commence with the special case treated at the beginning, with the same notations; \(g\) is measurable for the measure \(\pi(\mu) = \nu\) if and only if \(g_{\mathrm{L}}\) is \(\nu'\) -measurable (§1, No.~5, Example); now, this is equivalent to saying that \(g_{\mathrm{L}} \circ \pi' = (g \circ \pi)_{\mathrm{K}}\) is \(\mu_{\mathrm{K}}\) -measurable (Ch. V, §6, No.~2, Prop. 3), and finally that \(g \circ \pi\) is \(\mu\) -measurable (§1, No.~5, Example). Let us now pass to the general case, with the same notations as in the proof of Prop. 4; \(f\) is \(\nu\) -measurable if and only if \(f\) is \(\nu_{\alpha}\) -measurable for every \(\alpha \in \mathrm{A}\) (§1, No.~7, Prop. 8), hence if and only if \(f \circ \pi\) is \(\mu_{\alpha}\) -measurable for every \(\alpha \in \mathbf{A}\) (the preceding special case) and finally if and only if \(f \circ \pi\) is \(\mu\) -measurable (§1, No.~7, Prop. 8).

COROLLARY. --- Let X and Y be two topological spaces, \(\pi\) a \(\mu\) -proper mapping of T into X, and \(\pi'\) a \(\pi(\mu)\) -proper mapping of X into Y. The mapping \(\pi'' = \pi' \circ \pi\) is then \(\mu\) -proper, and \(\pi''(\mu) = \pi'(\pi(\mu))\) (`transitivity of images of measures').

For, \(\pi''\) is \(\mu\) -measurable (Prop. 5). Set \(\mu' = \pi(\mu)\) ; the image encumbrance \(\pi'(\mu'^{\bullet}) = \pi'(\pi(\mu^{\bullet}))\) is obviously equal to \(\pi''(\mu^{\bullet})\) . Since it is locally bounded, \(\pi''\) is \(\mu\) -proper. The measures \(\pi''(\mu)\) and \(\pi'(\mu')\) then have the same essential upper integral, hence are equal.

PROPOSITION 6. --- Let \(\pi\) be a \(\mu\) -proper mapping of T into a topological space X, and let B be a \(\pi(\mu)\) -measurable subset of X. Set A = \(\overline{\pi}^{-1}(B)\) , and denote by \(\pi'\) the mapping of A into B that coincides with \(\pi\) on A. The set A is then \(\mu\) -measurable, \(\pi_A\) and \(\pi'\) are \(\mu_A\) -proper, and

\[
\left(\pi (\mu)\right) _ {\mathrm{B}} = \left(\pi_ {\mathrm{A}} (\mu_ {\mathrm{A}})\right) _ {\mathrm{B}} = \pi^ {\prime} (\mu_ {\mathrm{A}}).\tag{2}
\]

The set A is \(\mu\) -measurable by Prop. 5 applied to \(\varphi_{B}\) ; the mapping \(\pi_{A}\) is clearly \(\mu_{A}\) -measurable by the definition of induced measures (No.~1), and it follows that \(\pi'\) is measurable. Let f be an element of \(\mathcal{F}_{+}(B)\) ; denoting by zero exponents the extensions by 0 in X and in T, we have

\[
\left(\pi (\mu) _ {\mathrm{B}}\right) ^ {\bullet} (f) = \pi (\mu) ^ {\bullet} \left(f ^ {0}\right) = \mu^ {\bullet} \left(f ^ {0} \circ \pi\right) = \mu^ {\bullet} \left(\left(f \circ \pi^ {\prime}\right) ^ {0}\right) = \mu_ {\mathrm{A}} ^ {\bullet} \left(f \circ \pi^ {\prime}\right),
\]

whence \(\left(\pi(\mu)_{\mathrm{B}}\right)^{\bullet} = \pi'(\mu_{\mathrm{A}}^{\bullet})\) . Since the encumbrance \(\left(\pi(\mu)_{\mathrm{B}}\right)^{\bullet}\) is locally bounded, the same is true of \(\pi'(\mu_{\mathrm{A}}^{\bullet})\) and therefore \(\pi'\) is \(\mu_{A}\) -proper. The measures \(\pi'(\mu_{\mathrm{A}})\) and \(\left(\pi(\mu)\right)_{\mathrm{B}}\) have the same essential upper integral, hence are equal. The other relation may be established in an analogous manner.

PROPOSITION 7. --- Let T be a subspace of a topological space X, and let i be the injection of T into X.

\begin{enumerate}
\def\labelenumi{\alph{enumi})}
\item
  If \(\mu\) is a measure on T, and if \(i\) is \(\mu\) -proper, then the measure \(i(\mu)\) is concentrated on T, and one has \((i(\mu))_{T} = \mu\) .
\item
  If \(\lambda\) is a measure on \(X\) such that \(T\) is \(\lambda\) -measurable, then \(i\) is \(\lambda_{T}\) -proper, and \(i(\lambda_T) = \varphi_T \cdot \lambda\) .
\item
  Set \(\nu = i(\mu)\) ; the relation \(\nu^{\bullet}(\mathbf{A}) = \mu^{\bullet}(\mathbf{A} \cap \mathbf{T})\) , applied to \(\mathbf{A} = \mathbf{X} - \mathbf{T}\) , shows that \(\nu\) is concentrated on \(\mathbf{T}\) . The relation \(\nu_{\mathrm{T}} = \mu\) is a special case of the relation (2), on taking \(\mathbf{B} = \mathbf{T} = \mathbf{A}\) .
\item
  Let \(f\) be a positive function defined on \(\mathbf{X}\) ; setting \(\mu = \lambda_{\mathrm{T}}\) , one has \(\mu^{\bullet}(f \circ i) = \lambda_{\mathrm{T}}^{\bullet}(f_{\mathrm{T}}) = \lambda^{\bullet}(f\varphi_{\mathrm{T}}) \leqslant \lambda^{\bullet}(f)\) (Prop. 1); it follows that \(i\) is \(\mu\) -proper. On the other hand, \(\mu^{\bullet}(f \circ i)\) (resp. \(\lambda^{\bullet}(f\varphi_{\mathrm{T}})\) ) is the essential upper integral of f with respect to \(i(\mu)\) (resp. \(\varphi_{T} \cdot \lambda\) ). These two measures are therefore equal.
\end{enumerate}

Remark 6). --- Let \(\pi\) be a \(\mu\) -proper mapping of T into a topological space X. One reduces the theory of integration with respect to the image measure \(\nu = \pi(\mu)\) to the theory treated in Ch. V, §6, in the following way. Let \((\mathrm{K}_{\alpha})_{\alpha \in \mathrm{A}}\) (resp. \((\mathrm{L}_{\beta})_{\beta \in \mathrm{B}}\) ) be a crushing of T (resp. of X) for \(\mu\) (resp. for \(\nu\) ), and set \(N = T - \bigcup_{\alpha \in A} K_{\alpha}\) , \(P = X - \bigcup_{\beta \in B} L_{\beta}\) . We can suppose that the restriction of \(\pi\) to each \(K_{\alpha}\) is continuous (§1, No.~8, Prop. 10). Let \(T', X'\) be the locally compact spaces constructed as in the Scholium of §1, No.~8 and let \(\mu'\) and \(\nu'\) be the measures on these spaces associated with \(\mu\) and \(\nu\) . The topology of \(T'\) being the sum of the topologies of the subspaces \(K_{\alpha}\) and the discrete topology on N, \(\pi\) is a continuous mapping of \(T'\) into X and the relation \(\mu'^{\bullet}(g \circ \pi) = \mu^{\bullet}(g \circ \pi) = \nu^{\bullet}(g)\) (for \(g \in \mathcal{F}_{+}(X)\) ) shows that \(\pi\) is \(\mu'\) -proper and that \(\pi(\mu') = \nu\) . On the other hand, the identity mapping \(i\) of X onto \(X'\) is \(\nu\) -proper, and \(i(\nu) = \nu'\) . It follows that \(\pi\) is a \(\mu'\) -proper mapping of \(T'\) into \(X'\) , and that the image of \(\mu'\) under \(\pi\) is \(\nu'\) (Cor. of Prop. 5). We leave to the reader the task of transcribing the results of Ch. V, §6.

\subsection{Lifting of measures}

PROPOSITION 8. --- Let T and X be two topological spaces, \(\pi\) a mapping of T into X.

\begin{enumerate}
\def\labelenumi{\alph{enumi})}
\item
  Let \(\nu\) be a bounded measure on X. In order that there exist a measure \(\mu\) on T such that \(\pi\) is \(\mu\) -proper and \(\pi(\mu) = \nu\) , it is necessary and sufficient that there exist, for every number \(\varepsilon > 0\) , a compact set \(\mathrm{K}_{\varepsilon} \subset \mathrm{T}\) such that the restriction of \(\pi\) to \(\mathrm{K}_{\varepsilon}\) is continuous and \(\nu^{\bullet}(\mathrm{X} - \pi(\mathrm{K}_{\varepsilon})) < \varepsilon\) .
\item
  Suppose that \(\pi\) is injective; let \(\mu\) and \(\mu'\) be two measures on T, such that \(\pi\) is proper for \(\mu\) and \(\mu'\) , and such that \(\pi(\mu) = \pi(\mu')\) . Then \(\mu = \mu'\) .
\end{enumerate}

The condition stated in \(a\) ) is necessary. For, if \(\pi\) is \(\mu\) -proper and \(\pi(\mu) = \nu\) , the relation \(\mu^{\bullet}(1) = \nu^{\bullet}(1) < +\infty\) implies that \(\mu\) is bounded. Prop. 2 of §1, No.~2, applied to the function 1, implies the existence of a compact subset K of T such that \(\mu^{\bullet}(\mathrm{T} - \mathrm{K}) < \varepsilon/2\) . Since \(\pi\) is \(\mu\) -measurable, there exists a compact set \(\mathrm{K}_{\varepsilon} \subset \mathrm{K}\) such that the restriction of \(\pi\) to \(\mathrm{K}_{\varepsilon}\) is continuous, and such that \(\mu^{\bullet}(\mathrm{K} - \mathrm{K}_{\varepsilon}) < \varepsilon/2\) . Then (No.~3, Prop. 4)

\[
\nu^ {\bullet} \big (\mathrm{X} - \pi (\mathrm{K} _ {\varepsilon}) \big) = \mu^ {\bullet} \big (\mathrm{T} - \pi^{-1} \big (\pi (\mathrm{K} _ {\varepsilon}) \big) \big) <   \varepsilon .
\]

To show that the condition is sufficient, we first treat a special case.

Lemma 1. --- Let U and V be two compact spaces, h a continuous mapping of U onto V. The mapping \(\lambda \mapsto h(\lambda)\) of \(\mathcal{M}_{+}(U)\) into \(\mathcal{M}_{+}(V)\) is then surjective.

For, let \(a\) be the linear mapping \(f \mapsto f \circ h\) of \(\mathcal{C}(\mathrm{V})\) into \(\mathcal{C}(\mathrm{U})\) ; since \(h\) is surjective, \(a\) is an isometry of \(\mathcal{C}(\mathrm{V})\) onto a subspace \(\mathrm{H}\) of \(\mathcal{C}(\mathrm{U})\) . Let \(\theta\) be a positive measure on \(\mathrm{V}\) ; then \(\theta \circ a^{-1}\) is a continuous linear form on \(\mathrm{H}\) , which is extendible to a linear form \(\eta\) on \(\mathcal{C}(\mathrm{U})\) with the same norm, by virtue of the Hahn--Banach theorem (TVS, II, §3, No.~2, Cor. 3 of Th. 1); \(\eta\) is then a measure on \(\mathrm{U}\) , and \(\theta(f) = \eta(f \circ h)\) for all \(f \in \mathcal{C}(\mathrm{V})\) , so that \(\theta = h(\eta)\) . Finally, \(\theta(1) = \| \theta \| = \| \eta \|\) , and \(\theta(1) = \eta(1)\) , so that \(\eta\) is positive (Ch. V, §5, No.~5, Prop. 9).

Let us now prove the sufficiency of the condition stated in \(a\) ). The condition implies the existence of a sequence \((\mathbf{K}_n)_{n \geqslant 1}\) of compact subsets of \(T\) , such that the restriction of \(\pi\) to each \(\mathbf{K}_n\) is continuous, and such that, for every \(n\) , \(\nu^{\bullet}(\mathrm{X} - \pi(\mathrm{K}_n)) < 1/n\) . The sequence \((\mathbf{K}_n)\) can be assumed to be increasing. Set \(\mathrm{L}_n = \pi(\mathrm{K}_n)\) and denote by \(\nu_n'\) the measure \(\varphi_{\mathrm{L}_n - \mathrm{L}_{n-1}} \cdot \nu_{\mathrm{L}_n}\) on \(\mathrm{L}_n\) , with the convention \(\mathrm{L}_0 = \emptyset\) .

The restriction \(\pi_{\mathrm{K}_n}\) being continuous, there exists a measure \(\mu_n'\) on \(\mathbf{K}_n\) such that \(\pi_{\mathrm{K}_n}(\mu_n') = \nu_n'\) (Lemma 1). Let \(\mu_n\) be the image of \(\mu_n'\) under the canonical injection of \(\mathbf{K}_n\) into \(\mathbf{T}\) , and let \(g\) be an element of \(\mathcal{F}_{+}(\mathbf{X})\) . Using successively the fact that \(\nu\) is concentrated on \(\bigcup_{n} \mathrm{L}_n\) , Prop. 4 of §1, No.~5; Prop. 2 of §1, No.~2; Prop. 4 of No.~3, and finally Prop. 7 of No.~3, we have

\[
\begin{array}{r l} \nu^ {\bullet} (g) & = \sum_ {n} \nu^ {\bullet} (\varphi_ {\mathrm{L} _ {n} - \mathrm{L} _ {n - 1}} g) = \sum_ {n} \nu_ {n} ^ {\prime   \bullet} (g _ {\mathrm{L} _ {n}}) = \sum_ {n} \mu_ {n} ^ {\prime   \bullet} (g _ {\mathrm{L} _ {n}} \circ \pi_ {\mathrm{K} _ {n}}) \\ & = \sum_ {n} \mu_ {n} ^ {\prime   \bullet} ((g \circ \pi) _ {\mathrm{K} _ {n}}) = \sum_ {n} \mu_ {n} ^ {\bullet} (g \circ \pi). \end{array}
\]

Taking \(g = 1\) in this formula, one sees that the family \((\mu_n)\) is summable and that its sum is a bounded measure \(\mu\) (§1, No.~7, Prop. 7). By Prop. 5 of No.~3, the mapping \(\pi\) is \(\mu_n\) -measurable for all \(n\) , because \(\pi_{\mathrm{K}_n}\) is continuous, hence \(\mu_n'\) -measurable; it follows that \(\pi\) is \(\mu\) -measurable (§1, No.~7, Prop. 8), hence \(\mu\) -proper since \(\mu\) is bounded. The above relations then prove that the measures \(\pi(\mu)\) and \(\nu\) have the same essential upper integral, hence are equal (§1, No.~2, Cor. of Prop. 2).

Finally, let us assume that \(\pi\) is injective, and let us prove \(b\) . Let \(f\) be an element of \(\mathcal{F}_{+}(\mathrm{T})\) ; since \(\pi\) is injective, there exists a function \(g \in \mathcal{F}_{+}(\mathrm{X})\) such that \(f = g \circ \pi\) and, setting \(\nu = \pi(\mu) = \pi(\mu')\) , by Prop. 4 of No.~3 we have

\[
\mu^ {\bullet} (f) = \mu^ {\bullet} (g \circ \pi) = \nu^ {\bullet} (g) = \mu^ {\prime \bullet} (g \circ \pi) = \mu^ {\prime \bullet} (f).
\]

The two measures \(\mu\) and \(\mu'\) thus have the same essential upper integral, which implies their equality (§1, No.~2, Cor. of Prop. 2).

Remark. --- Suppose that \(\pi\) is injective. Let \(\theta\) be a complex measure such that \(\pi\) is \(\theta\) -proper and \(\pi(\theta)=0\) ; then \(\theta=0\) . Indeed, by separating \(\theta\) into its real and imaginary parts, one can reduce to the case that \(\theta\) is real. We then have \(\pi(\theta^{+})=\pi(\theta^{-})\) , therefore \(\theta^{+}=\theta^{-}\) (Prop. 8), and finally \(\theta=0\) .

Here is an important case where condition \(a)\) of Prop. 8 is always satisfied.

PROPOSITION 9. --- Let T be a Souslin space (GT, IX, §6, No.~2, Def. 2), X a Hausdorff space, \(\pi\) a continuous mapping of T onto X, and \(\nu\) a bounded measure on X. Then there exists a bounded measure \(\mu\) on T such that \(\pi(\mu) = \nu\) .

The hypotheses obviously imply that \(\mathbf{X}\) is a Souslin space.

Let us consider the set function \(c: \mathrm{A} \mapsto \nu^{\bullet}(\pi(\mathrm{A}))\) on \(\mathfrak{P}(\mathrm{T})\) . The relation \(\mathrm{A} \subset \mathrm{B}\) implies \(c(\mathrm{A}) \leqslant c(\mathrm{B})\) ; if \((\mathrm{A}_n)\) is an increasing sequence of subsets of \(\mathrm{T}\) , and if \(\mathrm{A} = \bigcup_{n \in \mathbf{N}} \mathrm{A}_n\) , then \(c(\mathrm{A}) = \sup_n c(\mathrm{A}_n)\) from the fact that \(\nu^{\bullet}\) is an encumbrance. Finally, let \(\mathrm{A} \subset \mathrm{T}\) and let \(\varepsilon\) be a number \(> 0\) ; choose an open subset \(\mathrm{G}\) of \(\mathrm{X}\) containing \(\pi(\mathrm{A})\) , such that \(\nu^{\bullet}(\mathrm{G}) \leqslant \nu^{\bullet}(\pi(\mathrm{A})) + \varepsilon\) (§1, No.~9, Prop. 13); the open subset \(\mathrm{H} = \overline{\pi}^{-1}(\mathrm{G})\) of \(\mathrm{T}\) contains \(\mathrm{A}\) , and \(c(\mathrm{H}) \leqslant c(\mathrm{A}) + \varepsilon\) . The function \(c\) is therefore a right-continuous capacity on \(\mathrm{T} (\mathrm{TG},\mathrm{IX},\S 6,\mathrm{No.~10},\mathrm{Def.~9})^{(1)}\) and the theorem on capacitability (loc. cit., Th. 6) implies the equality \(c(\mathrm{T}) = \sup_{\mathrm{K}}c(\mathrm{K})\) , where \(K\) runs over the set of compact subsets of \(T\) . Prop. 8 then implies the existence of the desired measure \(\mu\) .

\subsection{Product of two measures}

Let S and T be two topological spaces, equipped respectively with two (positive) premeasures \(\lambda\) and \(\mu\) , and let X be the product space \(S \times T\) . Let K be a compact subset of X; let us denote by A and B the projections of K on S and T respectively, and set

\[
\nu_ {\mathrm{K}} = \left(\lambda_ {\mathrm{A}} \otimes \mu_ {\mathrm{B}}\right) _ {\mathrm{K}}.\tag{3}
\]

We thus define a premeasure on X. For, let L be a compact subset of X containing K, and let C and D be its two projections; then \(A \subset C\) , \(B \subset D\) , consequently, using the transitivity of induced measures, and Prop. 12 of Ch. V, §8, No.~5, we have

\[
\begin{array}{r l} & {(\nu_ {\mathrm{L}}) _ {\mathrm{K}} = \big ((\lambda_ {\mathrm{C}} \otimes \mu_ {\mathrm{D}}) _ {\mathrm{L}} \big) _ {\mathrm{K}} = (\lambda_ {\mathrm{C}} \otimes \mu_ {\mathrm{D}}) _ {\mathrm{K}}} \\ & {\qquad = \big ((\lambda_ {\mathrm{C}} \otimes \mu_ {\mathrm{D}}) _ {\mathrm{A} \times \mathrm{B}} \big) _ {\mathrm{K}} = (\lambda_ {\mathrm{A}} \otimes \mu_ {\mathrm{B}}) _ {\mathrm{K}} = \nu_ {\mathrm{K}}.} \end{array}
\]

DEFINITION 6. --- The premeasure \(\nu\) defined by (3) is called the product premeasure of \(\lambda\) and \(\mu\) , and is denoted \(\lambda \otimes \mu\) .

This definition obviously extends to the case that \(\lambda\) and \(\mu\) are complex premeasures, and one then has \(|\lambda\otimes\mu|=|\lambda|\otimes|\mu|\) (Ch. III, §4, No.~2, Prop. 3 and Ch. IV, §5, No.~7, Lemma 3).

We conserve the notations of Ch. III, §4 and Ch. V, §8 relative to products of measures and to iterated integrals. In particular, if \(f\) and \(g\) are two functions defined respectively on S and T, with values in \(\overline{\mathbf{R}}_+\) or in \(\mathbf{C}\) , the function \((s,t)\mapsto f(s)g(t)\) on \(S\times T\) will be denoted \(f\otimes g\) .

PROPOSITION 10. --- Let \(\nu\) be the product premeasure of \(\lambda\) and \(\mu\) ; for every function \(f \in \mathcal{F}_{+}(\mathrm{S})\) and every function \(g \in \mathcal{F}_{+}(\mathrm{T})\) ,

\[
\nu^ {\bullet} (f \otimes g) = \lambda^ {\bullet} (f) \mu^ {\bullet} (g).\tag{4}
\]

The premeasure \(\nu\) is the only premeasure on \(S \times T\) that satisfies (4).

As K (resp. L) runs over the set of compact subsets of S (resp. of T), we have

\[
\begin{array}{l} \nu^ {\bullet} (f \otimes g) = \sup _ {K, L} \nu_ {K \times L} ^ {\bullet} \big ((f \otimes g) _ {K \times L} \big) = \sup _ {K, L} (\lambda_ {K} \otimes \mu_ {L}) ^ {\bullet} (f _ {K} \otimes g _ {L}) \\ \qquad = \sup _ {K, L} \lambda_ {K} ^ {\bullet} (f _ {K}) \cdot \mu_ {L} ^ {\bullet} (g _ {L}) = \big (\sup _ {K} \lambda_ {K} ^ {\bullet} (f _ {K}) \big) \big (\sup _ {L} \mu_ {L} ^ {\bullet} (g _ {L}) \big) \\ \qquad = \lambda^ {\bullet} (f) \mu^ {\bullet} (g) \end{array}
\]

by Prop. 8 of Ch. V, §8, No.~3.

Let \(\eta\) be a second premeasure on \(S \times T\) satisfying (4), and let \(K\) and \(L\) be compact subsets of \(S\) and \(T\) respectively, \(f\) and \(g\) elements of \(\mathcal{F}_{+}(K)\) and \(\mathcal{F}_{+}(L)\) respectively. One has the relation \((f \otimes g)^{0} = f^{0} \otimes g^{0}\) between the extensions by 0, therefore (§1, No.~2, Prop. 2)

\[
\begin{array}{r l} & {\eta_ {\mathrm{K} \times \mathrm{L}} ^ {\bullet} (f \otimes g) = \eta^ {\bullet} \big ((f \otimes g) ^ {0} \big) = \eta^ {\bullet} (f ^ {0} \otimes g ^ {0})} \\ & {\qquad = \lambda^ {\bullet} (f ^ {0}) \mu^ {\bullet} (g ^ {0}) = \lambda_ {\mathrm{K}} ^ {\bullet} (f) \mu_ {\mathrm{L}} ^ {\bullet} (g).} \end{array}
\]

In particular, if one takes \(f \in \mathcal{K}_{+}(\mathrm{K})\) , \(g \in \mathcal{K}_{+}(\mathrm{L})\) , one sees that \(\eta_{\mathrm{K} \times \mathrm{L}}\) has the characteristic property of the product measure \(\lambda_{\mathrm{K}} \otimes \mu_{\mathrm{L}}\) (Ch. III, §4, No.~1, Th. 1). Therefore \(\eta_{\mathrm{K} \times \mathrm{L}} = \nu_{\mathrm{K} \times \mathrm{L}}\) ; since every compact subset of \(S \times T\) is contained in a set of the form \(K \times L\) , the transitivity of induced measures implies that \(\eta = \nu\) .

COROLLARY 1. --- If \(\lambda\) and \(\mu\) are measures, then \(\nu\) is a measure.

For, let \(x = (s, t)\) be a point of X, and let U and V be neighborhoods of s, t respectively, such that \(\lambda^{\bullet}(\mathrm{U}) < +\infty\) , \(\mu^{\bullet}(\mathrm{V}) < +\infty\) ; the set \(U \times V\) is a neighborhood of x, and \(\nu^{\bullet}(\mathrm{U} \times \mathrm{V}) = \lambda^{\bullet}(\mathrm{U}) \mu^{\bullet}(\mathrm{V}) < +\infty\) by (4); the encumbrance \(\nu^{\bullet}\) is thus locally bounded, and the premeasure \(\nu\) is a measure.

This result extends at once to complex measures.

COROLLARY 2. --- If A is a subset of S locally negligible for \(\lambda\) , then \(A \times T\) is locally \(\nu\) -negligible.

COROLLARY 3. --- Suppose that \(\lambda\) (resp. \(\mu\) ) is the sum of a summable family \((\lambda_{\alpha})_{\alpha \in A}\) (resp. \((\mu_{\beta})_{\beta \in B}\) ) of measures on S (resp. T). The family \((\lambda_{\alpha} \otimes \mu_{\beta})_{(\alpha, \beta) \in A \times B}\) is then summable, and its sum is \(\lambda \otimes \mu\) .

For, let \(p\) be the encumbrance \(\sum_{\alpha, \beta} (\lambda_{\alpha} \otimes \mu_{\beta})^{\bullet}\) ; if \(f \in \mathcal{F}_{+}(S)\) and \(g \in \mathcal{F}_{+}(T)\) , then obviously \(p(f \otimes g) = \lambda^{\bullet}(f)\mu^{\bullet}(g)\) . The proof of Cor. 1 then shows that \(p\) is locally bounded, so that the family \((\lambda_{\alpha} \otimes \mu_{\beta})\) is summable (§1, No.~7, Prop. 7). Its sum \(\eta\) then satisfies \(\eta^{\bullet} = p\) (§1, No.~7, Prop. 7), and Prop. 10 implies \(\eta = \nu\) .

\subsection{Integration with respect to the product of two measures}

Throughout this No., \(\lambda\) and \(\mu\) denote measures on S and T, respectively, and \(\nu\) denotes the product measure \(\lambda \otimes \mu\) on \(S \times T\) . In addition, if \(f\) is a positive function on \(S \times T\) , for every \(s \in S\) we denote by \(f_s\) the function \(t \mapsto f(s,t)\) on T, and by \(I_f\) the function \(s \mapsto \mu^\bullet(f_s)\) on S.

Lemma 2. --- Let \(f\) be a \(\nu\) -measurable positive function on \(S \times T\) ; for every compact subset \(L\) of \(T\) , let \(I_f^L\) be the function \(s \mapsto \mu^\bullet(f_s \varphi_L)\) on \(S\) . Then the function \(I_f^L\) is \(\lambda\) -measurable, and

\[
\mathrm{I} _ {f} = \sup _ {\mathrm{L}} \mathrm{I} _ {f} ^ {\mathrm{L}}\tag{5}
\]

\[
\nu^ {\bullet} (f) = \sup _ {L} \lambda^ {\bullet} (I _ {f} ^ {L}),\tag{6}
\]

where L runs over the set of compact subsets of T.

We first note that the inclusion \(L \subset L'\) implies \(I_{f}^{L} \leqslant I_{f}^{L'}\) ; on the other hand, \(\mathrm{I}_{f}^{\mathrm{L}}(s) = \mu_{\mathrm{L}}^{\bullet}((f_{s})_{\mathrm{L}})\) for all \(s \in S\) . Formula (5) is therefore an immediate consequence of the definition of the encumbrance \(\mu^{\bullet}\) given in §1, No.~2. If K is a compact subset of S, and L a compact subset of T, then \(\nu_{\mathrm{K} \times \mathrm{L}} = \lambda_{\mathrm{K}} \otimes \mu_{\mathrm{L}}\) by construction, and Prop. 7 of Ch. V, §8, No.~3 implies the relation

\[
\nu^ {\bullet} (f \varphi_ {\mathrm{K} \times \mathrm{L}}) = \lambda_ {\mathrm{K}} ^ {\bullet} \big ((\mathrm{I} _ {f} ^ {\mathrm{L}}) _ {\mathrm{K}} \big).\tag{7}
\]

Moreover, every compact subset of \(S \times T\) is contained in a compact set of the form \(K \times L\) ; passing to the upper envelope over all K and L in the preceding formula, we therefore obtain

\[
\nu^ {\bullet} (f) = \sup _ {L} \sup _ {K} \lambda_ {K} ^ {\bullet} \big ((I _ {f} ^ {L}) _ {K} \big) = \sup _ {L} \lambda^ {\bullet} (I _ {f} ^ {L}),\tag{8}
\]

namely (6).

Finally, Prop. 7 of Ch. V, §8, No.~3 implies that the restriction of \(I_{f}^{L}\) to every compact subset K of S is \(\lambda_{K}\) -measurable; this is equivalent to saying that \(I_{f}^{L}\) is \(\lambda\) -measurable.

PROPOSITION 11. --- Let \(f\) be a lower semi-continuous function \(\geqslant 0\) defined on \(X = S \times T\) .

\begin{enumerate}
\def\labelenumi{\alph{enumi})}
\item
  The function \(f_{s}: t \mapsto f(s, t)\) is lower semi-continuous on \(\mathbf{T}\) for every \(s \in S\) .
\item
  The function \(I_{f}: s \mapsto \int^{\bullet} f(s, t) d\mu(t)\) is lower semi-continuous on S, and
\end{enumerate}

\[
\iint_ {\mathrm{X}} ^ {\bullet} f (s, t) d \nu (s, t) = \int_ {\mathrm{S}} ^ {\bullet} d \lambda (s) \int_ {\mathrm{T}} ^ {\bullet} f (s, t) d \mu (t).\tag{9}
\]

The property \(a\) ) is obvious, since the mapping \(t \mapsto f(s,t)\) of T into \(\overline{\mathbf{R}}\) is the composition of \(f\) with the continuous mapping \(t \mapsto (s,t)\) of T into X. To establish \(b\) ), we will make use of a lemma:

Lemma 3. --- Let X be a topological space (Hausdorff or not), f a lower semi-continuous function \(\geqslant 0\) defined on X; then f is the limit of an increasing sequence \((f_{n})_{n\in\mathbf{N}}\) of lower semi-continuous functions on X, such that each function \(f_{n}\) is a linear combination, with positive coefficients, of characteristic functions of open sets.

Given two integers \(k \geqslant 1\) and \(n \geqslant 1\) , let us denote by \(J_{kn}\) the characteristic function of the interval \(]k/2^n, +\infty]\) of \(\overline{\mathbf{R}}\) . For every \(x \in \overline{\mathbf{R}}_+\) , set \(u_n(x) = 2^{-n} \sum_{k=1}^{n \cdot 2^n} J_{kn}(x)\) ; it is immediate that the sequence \((u_n(x))_{n \geqslant 1}\) is increasing and admits \(x\) as limit. The sequence of functions \(f_n = u_n \circ f\) is therefore increasing and converges to \(f\) , and one has \(f_{n} = 2^{-n}\sum_{k = 1}^{n\cdot 2^{n}}\varphi_{\mathrm{U}(k,n)}\) ,

where \(\mathrm{U}(k,n)\) is the open set \(\overline{f} (]k / 2^n, + \infty ])\) of X.

Let us pass to the proof of \(b\) ). The function \(I_f\) being the upper envelope of the increasing directed family of functions \(I_f^L\) , where \(L\) runs over the set of compact subsets of \(T\) (Lemma 2), it will suffice to show that the functions \(I_f^L\) are lower semi-continuous; the formula (9) may then be deduced from (6) by passing to the upper envelope over \(L\) (§1, No.~6, Prop. 5).

Thus let \(\mathcal{H}\) be the set of positive lower semi-continuous functions \(f\) on \(\mathrm{S} \times \mathrm{T}\) such that \(\mathrm{I}_f^{\mathrm{L}}\) is lower semi-continuous for every compact subset \(\mathrm{L}\) of \(\mathrm{T}\) . By Prop. 5 of §1, No.~6, the supremum of every increasing directed set of elements of \(\mathcal{H}\) belongs to \(\mathcal{H}\) . By Lemma 3, it will therefore suffice to prove that the characteristic function of an open set \(\mathrm{W}\) of \(\mathrm{S} \times \mathrm{T}\) belongs to \(\mathcal{H}\) . Moreover, by the definition of the product topology on \(\mathrm{S} \times \mathrm{T}\) , the open set \(\mathrm{W}\) is the union of an increasing directed family \((\mathrm{W}_{\alpha})_{\alpha \in \mathrm{A}}\) of open sets of the form

\[
\mathrm{W} = \bigcup_ {1 \leqslant i \leqslant n} \left(\mathrm{U} _ {i} \times \mathrm{V} _ {i}\right),
\]

where the \(U_{i}\) are open in S and the \(V_{i}\) are open in T; by the remarks made above, it will suffice to show that the characteristic function of such an open set belongs to H. Let then \(s \in S\) , and let U be the intersection of the family (possibly empty) formed by the open sets \(U_{i}\) containing s; one sees immediately that \(\varphi_{\mathrm{W}}(s,t) \leqslant \varphi_{\mathrm{W}}(s',t)\) for all \(s' \in U\) and \(t \in T\) , whence, by integration, \(\mathrm{I}_{\varphi_{\mathrm{W}}}^{\mathrm{L}}(s) \leqslant \mathrm{I}_{\varphi_{\mathrm{W}}}^{\mathrm{L}}(s')\) for all \(s' \in U\) . Consequently \(I_{\varphi_{w}}^{L}\) is lower semi-continuous, and the proposition is established.

COROLLARY 1. --- Let \(f\) be a positive numerical function defined on \(\mathbf{X} = \mathbf{S} \times \mathbf{T}\) ; then

\[
\iint_ {\mathrm{X}} ^ {*} f (s, t) d \nu (s, t) \geqslant \int_ {\mathrm{S}} ^ {*} d \lambda (s) \int_ {\mathrm{T}} ^ {*} f (s, t) d \mu (t).\tag{10}
\]

For, let \(g\) be a lower semi-continuous function on \(\mathbf{X}\) such that \(g \geqslant f\) ; by Prop. 11,

\[
\begin{array}{r l} \iint_ {\mathrm{X}} ^ {*} g (s, t) d \nu (s, t) & = \iint^ {\bullet} g (s, t) d \nu (s, t) = \int^ {\bullet} d \lambda (s) \int^ {\bullet} g (s, t) d \mu (t) \\ & = \int^ {*} d \lambda (s) \int^ {*} g (s, t) d \mu (t) \geqslant \int^ {*} d \lambda (s) \int^ {*} f (s, t) d \mu (t). \end{array}
\]

The inequality (10) is obtained by passing to the lower envelope over \(g\) .

COROLLARY 2. --- Let f be a numerical function defined on S × T and ν-negligible. Then the function \(f_{s}: t \mapsto f(s, t)\) is μ-negligible for λ-almost every \(s \in S\) .

PROPOSITION 12. --- Let f be a ν-measurable positive function defined on X = S × T. Assume that f is ν-moderated (resp. that μ is moderated). Then:

\begin{enumerate}
\def\labelenumi{\alph{enumi})}
\item
  The set N of \(s \in S\) such that the function \(f_{s}: t \mapsto f(s, t)\) is not \(\mu\) -measurable is negligible (resp. locally negligible) for \(\lambda\) .
\item
  The mapping \(s \mapsto \int^{\bullet} f(s, t) \, d\mu(t)\) is \(\lambda\) -measurable, and
\end{enumerate}

\[
\iint_ {\mathrm{X}} ^ {\bullet} f (s, t) d \nu (s, t) = \int_ {\mathrm{S}} ^ {\bullet} d \lambda (s) \int_ {\mathrm{T}} ^ {\bullet} f (s, t) d \mu (t).\tag{11}
\]

We begin by establishing \(b)\) when \(f\) is \(\nu\) -moderated. By Lemma 2, this part of the statement is valid when there exists a compact subset \(L\) of \(T\) such that \(f\) is zero outside \(S \times L\) ; for, in this case \(I_f = I_f^{L'}\) for every compact subset \(L'\) of \(T\) containing \(L\) , and formula (11) reduces to (6). In particular, \(b)\) is established for a function \(f\) that is zero outside a compact subset of \(S \times T\) . On the other hand, Cor. 1 of Prop. 11 implies that \(b)\) is true when \(f\) is \(\nu\) -negligible. Since every \(\nu\) -moderated function is the sum of a \(\nu\) -negligible function and a sequence of functions with compact support (§1, No.~9, Cor. 3 of Prop. 14), the assertion \(b)\) is true when \(f\) is \(\nu\) -moderated.

Similarly, the assertion \(b)\) is obvious when \(\mu\) is carried by a compact subset \(L\) of \(T\) (Lemma 2). Suppose that \(\mu\) is moderated; then there exists a sequence \((\mu_n)_{n \in \mathbb{N}}\) of measures on \(T\) with compact support, such that \(\mu = \sum_{n} \mu_n\) (§1, No.~9, Cor. 5 of Prop. 14), whence \(\nu = \sum_{n} \lambda \otimes \mu_n\) (No.~5, Cor. 3 of Prop. 10). The assertion \(b)\) , being valid for each of the measures \(\nu_n = \lambda \otimes \mu_n\) , is also valid for \(\nu = \sum_{n} \nu_n\) .

Let us prove \(a\) ); denote by N the set of \(s \in S\) such that \(f_s\) is not \(\mu\) -measurable; for every compact subset L of T, similarly denote by \(\mathbf{N}_{\mathrm{L}}\) the set of \(s \in S\) such that \(f_s \varphi_{\mathrm{L}}\) is not \(\mu\) -measurable. If K and L are compact sets in S and T respectively, then \(f_{\mathrm{K} \times \mathrm{L}}\) is measurable with respect to the measure \(\nu_{\mathrm{K} \times \mathrm{L}} = \lambda_{\mathrm{K}} \otimes \mu_{\mathrm{L}}\) , and Prop. 2 of Ch. V, §8, No.~2 shows that the set \(\mathbf{N}_{\mathrm{L}}\) is locally negligible for \(\lambda_{\mathrm{K}}\) ; since K is arbitrary, \(\mathbf{N}_{\mathrm{L}}\) is locally \(\lambda\) -negligible.

Suppose that f is zero outside a compact set of the form \(K \times L\) ; then \(N = N_{L}\) , and N is contained in K; it follows that N is \(\lambda\) -negligible. Similarly, if f is \(\nu\) -negligible, Cor. 2 of Prop. 11 implies that N is \(\lambda\) -negligible. The case that f is \(\nu\) -moderated can then be treated as above, on combining the preceding two cases.

Suppose that \(\mu\) is carried by a compact subset L of T; then \(N = N_{L}\) again, therefore N is locally \(\lambda\) -negligible. Since every moderated measure is the sum of a sequence of measures with compact support ( \(\S1\) , No.~9, Cor. 5 of Prop. 14), this result extends at once to the case that \(\mu\) is moderated, on using Prop. 8 of \(\S1\) , No.~7.

Remark. --- Let \((\mathrm{K}_{\alpha})_{\alpha\in\mathrm{A}}\) be a crushing of S for \(\lambda\) and let \(M = S - \bigcup_{\alpha\in A} K_{\alpha}\) ; define in an analogous way \((\mathrm{L}_{\beta})_{\beta\in\mathrm{B}}\) and N for the measure \(\mu\) on T. We denote by \(S'\) the locally compact space that is the sum of the subspaces \(K_{\alpha}\) of S and the discrete space M; the space \(T'\) is defined analogously, and we set \(X' = S' \times T'\) . The locally compact space \(X'\) is the sum of the family \((\mathrm{K}_{\alpha} \times \mathrm{L}_{\beta})_{(\alpha,\beta)\in A\times B}\) of compact subspaces of X and the subspace \(P = (M \times T) \cup (S \times N)\) that is a locally \(\nu\) -negligible subset of X (one observes that P is in general not a discrete space). We saw in the Scholium of §1, No.~8 that there exists a measure \(\lambda'\) on \(S'\) such that the measurable functions, the essential upper integral of positive functions, the essentially integrable functions and their integrals, are the same for \(\lambda\) and \(\lambda'\) . We associate the measure \(\mu'\) on \(T'\) with \(\mu\) , and the measure \(\nu'\) on \(X'\) with \(\nu\) , in conformity with the cited Scholium; one sees immediately that \(\nu'\bullet(f \otimes g) = \lambda'\bullet(f)\mu'\bullet(g)\) for \(f \in \mathcal{F}_{+}(S)\) and \(g \in \mathcal{F}_{+}(T)\) ; therefore \(\nu' = \lambda' \otimes \mu'\) by Prop. 10 of No.~5. Since the topology of \(X'\) is finer than that of X, every \(\nu\) -moderated function is \(\nu'\) -moderated. This procedure permits extending without new proof the Lebesgue--Fubini theorem (Ch. V, §8, No.~4, Th. 1) to the present situation.

\subsection{A result on the disintegration of measures}

PROPOSITION 13. --- Let X be a topological space, \(\nu\) a moderated measure on X, p a \(\nu\) -proper mapping of X into a topological space T, and \(\mu = p(\nu)\) . Assume that every compact subspace of X is metrizable. Then there exists a mapping \(t \mapsto \lambda_{t}\) of T into \(\mathcal{M}_{+}(X)\) having the following properties:

\begin{enumerate}
\def\labelenumi{\alph{enumi})}
\item
  for every \(t \in \mathbf{T}\) , the measure \(\lambda_t\) is carried by \(\overline{p}^{-1}(t)\) ;
\item
  for every universally measurable \(^{(1)}\) positive function f on X, the function \(t \mapsto \lambda_{t}^{\bullet}(f)\) is universally measurable on T and
\end{enumerate}

\[
\int_ {\mathrm{X}} ^ {\bullet} f (x) d \nu (x) = \int_ {\mathrm{T}} ^ {\bullet} d \mu (t) \int_ {\mathrm{X}} ^ {\bullet} f (x) d \lambda_ {t} (x);\tag{12}
\]

\begin{enumerate}
\def\labelenumi{\alph{enumi})}
\setcounter{enumi}{2}
\tightlist
\item
  the set of \(t \in \mathrm{T}\) such that \(\lambda_t(1) \neq 1\) is locally \(\mu\) -negligible.
\end{enumerate}

Moreover, if \(t \mapsto \lambda_t'\) is a mapping of \(\mathbf{T}\) into \(\mathcal{M}_+(\mathbf{X})\) satisfying the conditions \(a)\) and \(b)\) , the set of \(t \in \mathbf{T}\) such that \(\lambda_t \neq \lambda_t'\) is locally \(\mu\) -negligible. We will need an auxiliary result:

Lemma 4. --- Let X be a topological space, \(\nu\) a measure on X, and \(f\) a \(\nu\) -measurable mapping of X into a topological space F (Hausdorff or not). There exists a universally measurable mapping \(f'\) of X into F, equal to \(f\) locally \(\nu\) -almost everywhere.

The proof is identical to that of Prop. 7 of Ch. V, §3, No.~4, on taking into account Prop. 10 of §1, No.~8.

Let us pass to the proof of Prop. 13.

\begin{enumerate}
\def\labelenumi{\Alph{enumi})}
\tightlist
\item
  Suppose that \(\mathbf{X}\) is compact and metrizable and that \(p\) is continuous and surjective:
\end{enumerate}

The space T is then compact and metrizable (GT, IX, §2, No.~10). By Th. 1 of Ch. VI, §3, No.~1, there exists a mapping H: \(t \mapsto \eta_t\) of T into \(\mathcal{M}_+(X)\) , vaguely \(\mu\) -measurable and scalarly essentially \(\mu\) -integrable, such that \(\nu = \int_{\mathrm{T}} \eta_t d\mu(t)\) and such that \(\eta_t\) has total mass 1 and is carried by \(\overline{p}^{-1}(t)\) for every \(t \in \mathrm{T}\) . Let \((\mathrm{S}_n)_{n \in \mathbf{N}}\) be a crushing of T for \(\mu\) , such that the restriction of H to each of the sets \(\mathrm{S}_n\) is continuous (§1, No.~8, Props. 10 and 11); we denote by \(\Lambda: t \mapsto \lambda_t\) the mapping of T into \(\mathcal{M}_+(X)\) that is equal to H on \(S = \bigcup_{n \in \mathbf{N}} S_n\) and to 0 on T - S. It is clear that \(\nu = \int_{\mathrm{T}}\lambda_t d\mu (t)\) and that \(\Lambda\) satisfies condition \(a)\) of the statement.

Let \(\theta\) be a measure on T; the mapping \(\Lambda\) is vaguely \(\theta\) -measurable and scalarly essentially \(\theta\) -integrable, hence also \(\theta\) -adequate (Ch. V, §3, No.~1, Prop. 2 b)). Let \(f\) be a universally measurable positive function on X; by Prop. 5 of Ch. V, §3, No.~2 applied to \(\int \lambda_t d\theta(t)\) , the mapping \(t \mapsto \lambda_t^{\bullet}(f)\) is \(\theta\) -measurable, hence universally measurable in view of the arbitrariness of \(\theta\) .

The formula (12) results from Prop. 5 of Ch. V, §3, No.~2.

\begin{enumerate}
\def\labelenumi{\Alph{enumi})}
\setcounter{enumi}{1}
\tightlist
\item
  Suppose that there exists a compact subset \(\mathbf{X}'\) of \(\mathbf{X}\) carrying the measure \(\nu\) and such that \(p_{\mathbf{X}'}\) is continuous:
\end{enumerate}

Set \(T' = p(X')\) , and \(p' = p_{X'}\) ; let us denote by \(\nu'\) the measure \(\nu_{X'}\) , and by \(\mu'\) the image measure \(p'(\nu')\) on \(T'\) . Since \(p'\) is continuous and surjective and since \(X'\) is compact and metrizable, by A) there exists a mapping \(\Lambda': t' \mapsto \lambda_{t'}'\) of \(T'\) into \(\mathcal{M}_{+}(X')\) satisfying the following conditions:

\(a^{\prime})\) for every \(t^\prime \in \mathrm{T}^\prime\) , the measure \(\lambda_{t'}'\) is carried by \(\mathrm{X}' \cap \overline{p}^{-1}(t')\) ;

\(b')\) for every universally measurable positive function \(f'\) on \(X'\) , the function \(t' \mapsto \lambda_{t'}^{\bullet}(f')\) is universally measurable on \(T'\) and

\[
\int_ {\mathrm{X} ^ {\prime}} ^ {\bullet} f ^ {\prime} (x ^ {\prime}) d \nu^ {\prime} (x ^ {\prime}) = \int_ {\mathrm{T} ^ {\prime}} ^ {\bullet} d \mu^ {\prime} (t ^ {\prime}) \int_ {\mathrm{X} ^ {\prime}} ^ {\bullet} f ^ {\prime} (x ^ {\prime}) d \lambda_ {t ^ {\prime}} ^ {\prime} (x ^ {\prime}).
\]

Let \(t \in \mathbf{T}\) ; if \(t\) belongs to \(\mathbf{T}'\) , let us denote by \(\lambda_t\) the image of \(\lambda_t'\) under the canonical injection of \(\mathbf{X}'\) into \(\mathbf{X}\) , and if \(t\) belongs to \(\mathbf{T} - \mathbf{T}'\) we set \(\lambda_t = 0\) . The reader will verify without difficulty that the mapping \(t \mapsto \lambda_t\) satisfies the conditions \(a)\) and \(b)\) of the statement.

\begin{enumerate}
\def\labelenumi{\Alph{enumi})}
\setcounter{enumi}{2}
\tightlist
\item
  Existence in the general case:
\end{enumerate}

The measure \(\nu\) on X being moderated, we may choose a covering \((\mathrm{U}_m)_{m\in \mathbf{N}}\) of X consisting of \(\nu\) -integrable open sets. Let in addition \((\mathrm{X}_n)_{n\in \mathbf{N}}\) be a \(\nu\) -crushing of X such that the restriction of \(p\) to each set \(\mathrm{X}_n\) is continuous (§1, No.~8, Props. 10 and 11); denote by \(\nu_n\) the measure \(\varphi_{\mathrm{X}_n}\cdot \nu\) on X and by \(\mu_n\) its image under \(p\) . By B) there exists, for every integer \(n\in \mathbf{N}\) , a mapping \(t\mapsto \alpha_t^n\) of T into \(\mathcal{M}_{+}(\mathrm{X})\) satisfying the following conditions:

\(a^{\prime \prime})\) The measure \(\alpha_{t}^{n}\) is carried by \(\overline{p}^{-1}(t)\) for all \(t\in \mathbf{T}\) .

\(b^{\prime\prime}\) If f is a universally measurable positive function on X, the positive function \(t \mapsto (\alpha_{t}^{n})^{\bullet}(f)\) on T is universally measurable and

\[
\int_ {\mathrm{X}} ^ {\bullet} f (x) d \nu_ {n} (x) = \int_ {\mathrm{T}} ^ {\bullet} d \mu_ {n} (t) \int_ {\mathrm{X}} ^ {\bullet} f (x) d \alpha_ {t} ^ {n} (x).\tag{13}
\]

One has \(\nu = \sum_{n\in \mathbf{N}}\nu_n\) and \(\mu = \sum_{n\in \mathbf{N}}\mu_n\) ; it follows immediately from Prop. 3 of No.~2 and the above Lemma 4 that there exists a sequence \((g_n)_{n\in \mathbf{N}}\) of universally measurable positive functions on T such that \(\mu_{n} = g_{n}\cdot \mu\) for all \(n\in \mathbf{N}\) and such that \(\sum_{n\in \mathbf{N}}g_n = 1\) . For every \(t\in \mathbf{T}\) , let us denote by \(\beta_t^n\) the measure \(g_{n}(t)\cdot \alpha_{t}^{n}\) on X and by \(q_{t}\) the encumbrance \(\sum_{n\in \mathbf{N}}(\beta_t^n)^{\bullet}\) on X. Let \(f\) be a universally measurable positive function on X; using Prop. 2 of No.~2 and summing over \(n\) in (13), we obtain

\[
\int_ {\mathrm{X}} ^ {\bullet} f (x) d \nu (x) = \int_ {\mathrm{T}} ^ {\bullet} q _ {t} (f) d \mu (t);\tag{14}
\]

it is moreover clear that the function \(t \mapsto q_t(f)\) on \(\mathbf{T}\) is universally measurable.

For every \(m \in \mathbf{N}\) , let \(\mathbf{E}_m\) be the set of \(t \in \mathbf{T}\) such that \(q_t(\mathbf{U}_m) = +\infty\) ; the set \(\mathbf{E}_m\) is universally measurable because this is true of the mapping \(t \mapsto q_t(\mathbf{U}_m)\) , and \(\mathbf{E}_m\) is locally \(\mu\) -negligible by the formula (14) applied to \(f = \varphi_{\mathbf{U}_m}\) , since \(\nu^\bullet(\mathbf{U}_m)\) is finite. The set \(\mathbf{E} = \bigcup_{m \in \mathbf{N}} \mathbf{E}_m\) is therefore universally measurable and locally \(\mu\) -negligible. We set \(\lambda_t = 0\) for \(t \in \mathbf{E}\) . Moreover, let \(t \in \mathbf{T} - \mathbf{E}\) ; the encumbrance \(q_t\) is locally bounded since the open sets \(\mathbf{U}_m\) cover \(\mathbf{X}\) and since \(q_t(\mathbf{U}_m)\) is finite; by Prop. 7 of §1, No.~7, there exists a measure \(\lambda_t\) on \(\mathbf{X}\) such that \(q_t = \lambda_t^\bullet\) and \(\lambda_t = \sum_{n \in \mathbf{N}} \beta_t^n\) . It is immediate that the mapping \(t \mapsto \lambda_t\) satisfies the conditions \(a)\) and \(b)\) of the statement.

D) Proof of c):

Let \(f\) be a universally measurable function on \(X\) that is positive and bounded; we are going to show that the universally measurable function \(h_f: t \mapsto \lambda_t^\bullet(f)\) on \(T\) is a density for the measure \(\mu_f = p(f \cdot \nu)\) with respect to \(\mu = p(\nu)\) . Let \(K\) be a compact subset of \(T\) and let \(A = \overline{p}^{-1}(K)\) . For every \(t \in T\) , the measure \(\lambda_t\) is carried by \(\overline{p}^{-1}(t)\) ; if \(t\) belongs to \(K\) then \(\overline{p}^{-1}(t) \subset A\) , whence \(\lambda_t^\bullet(f\varphi_A) = \lambda_t^\bullet(f)\) ; on the other hand, if \(t\) belongs to \(T - K\) then \(\overline{p}^{-1}(t) \subset X - A\) , whence \(\lambda_t^\bullet(f\varphi_A) = 0\) . Applying the formula (12) to \(f \cdot \varphi_A\) , \(^{(1)}\) we obtain

\[
\mu_ {f} (\mathbf {K}) = \int_ {\mathrm{A}} ^ {\bullet} f (x) d \nu (x) = \int_ {\mathrm{K}} ^ {\bullet} d \mu (t) \int_ {\mathrm{X}} ^ {\bullet} f (x) d \lambda_ {t} (x) = \int_ {\mathrm{K}} ^ {\bullet} h _ {f} (t) d \mu (t),
\]

which establishes the relation \(\mu_f = h_f \cdot \mu\) .

Letting \(f = 1\) , one sees that the function \(h_1: t \mapsto \| \lambda_t \|\) is a density of the measure \(\mu_1 = \mu\) with respect to \(\mu\) , hence is equal to 1 locally \(\mu\) -almost everywhere in T.

E) Uniqueness:

Let \(t \mapsto \lambda_t^i\) (for \(i = 1,2\) ) be two mappings of \(T\) into \(\mathcal{M}_+(X)\) satisfying the conditions \(a)\) and \(b)\) of the statement. As in C), choose a \(\mu\) -crushing \((X_n)_{n \in \mathbb{N}}\) of \(X\) such that \(p_{X_n}\) is continuous for every \(n \in \mathbb{N}\) , and set \(N = X - \bigcup_{n \in \mathbb{N}} X_n\) . For every integer \(n \in \mathbb{N}\) , choose a countable set \(D_n\) of positive functions on \(X\) , zero outside \(X_n\) , whose restrictions to \(X_n\) form a dense set in the normed space \(\mathcal{C}(X_n)\) (apply Th. 1 of GT, \(X, \S 3\) , No.~3 to the metrizable compact space \(X_n\) ). We set \(D = \bigcup_{n \in \mathbb{N}} D_n\) .

Let \(f \in \mathbf{D}\) ; by D), the functions \(t \mapsto (\lambda_t^1)^{\bullet}(f)\) and \(t \mapsto (\lambda_t^2)^{\bullet}(f)\) are densities of the measure \(\mu_f\) with respect to \(\mu\) , and so there exists a locally \(\mu\) -negligible set \(\mathbf{E}_f\) in T such that \((\lambda_t^1)^{\bullet}(f) = (\lambda_t^2)^{\bullet}(f)\) for \(t \in \mathbf{T} - \mathbf{E}_f\) . Moreover, by (12), the set \(\mathbf{F}_i\) of \(t \in \mathbf{T}\) such that \((\lambda_t^i)^{\bullet}(\mathbf{N}) \neq 0\) is locally \(\mu\) -negligible for \(i = 1, 2\) . Since D is countable, the set \(G = \left( \bigcup_{f \in D} E_f \right) \cup F_1 \cup F_2\) is locally \(\mu\) -negligible; for \(t \in T - G\) , we have \((\lambda_t^1)^{\bullet}(\mathbf{N}) = (\lambda_t^2)^{\bullet}(\mathbf{N}) = 0\) and

\((\lambda_t^1)_{\mathbf{X}_n} = (\lambda_t^2)_{\mathbf{X}_n}\) , whence \(\lambda_t^1 = \lambda_t^2\) by Prop. 9 of §1, No.~8.

Q.E.D.

Remarks. --- 1) If X is a Souslin space, then every compact subspace of X is a Souslin space, hence is metrizable (TG, IX, Appendix I, Cor. 2 of Prop. 3), \(^{(2)}\) and every measure on X is moderated (§1, No.~9, Remark 1). By Prop. 13, every measure \(\nu\) on X therefore admits a disintegration with respect to every \(\nu\) -proper mapping.

\begin{enumerate}
\def\labelenumi{\arabic{enumi})}
\setcounter{enumi}{1}
\tightlist
\item
  With the notations of Prop. 13, let \(f\) be a positive \(\nu\) -measurable function. One can prove, as in Ch. V, §3, No.~2, Prop. 5, that the set of \(t \in \mathbf{T}\) such that \(f\) is not \(\lambda_t\) -measurable is locally \(\mu\) -negligible, that \(t \mapsto \lambda_t^\bullet(f)\) is \(\mu\) -measurable, and that the relation (12) again holds.
\end{enumerate}

\section{MEASURES AND ADDITIVE SET FUNCTIONS}

In this section, we shall denote by \(\mathfrak{K}(\mathrm{T})\) and \(\mathfrak{B}(\mathrm{T})\) , respectively, the set of compact subsets of a Hausdorff topological space \(\mathrm{T}\) and the Borel tribe of \(\mathrm{T}\) .

\subsection{Measures and additive set functions of compact sets}

THEOREM 1. --- Let T be a topological space, and I a mapping of \(\mathfrak{K}(\mathrm{T})\) into \(\mathbf{R}_{+}\) . In order that there exist a measure \(\mu\) on T such that \(\mathrm{I}(\mathrm{K}) = \mu^{\bullet}(\mathrm{K})\) for all \(\mathrm{K} \in \mathfrak{K}(\mathrm{T})\) , it is necessary and sufficient that I satisfy the following conditions:

\begin{enumerate}
\def\labelenumi{\arabic{enumi})}
\item
  If K and L are compact subsets of T such that \(K \subset L\) , then \(I(K) \leqslant I(L)\) (`I is increasing').
\item
  If K and L are compact subsets of T, then I(K ∪ L) ≤ I(K) + I(L).
\item
  If K and L are disjoint compact subsets of T, then \(\mathrm{I}(\mathrm{K}\cup\mathrm{L})=\mathrm{I}(\mathrm{K})+\mathrm{I}(\mathrm{L})\) (`I is additive').
\item
  For every decreasing directed family \((\mathrm{K}_{\alpha})_{\alpha \in \mathrm{A}}\) of compact subsets of T, one has \(\mathrm{I}\big(\bigcap_{\alpha \in \mathrm{A}}\mathrm{K}_{\alpha}\big) = \inf_{\alpha \in \mathrm{A}}\mathrm{I}(\mathrm{K}_{\alpha})\) .
\item
  For every \(x \in \mathbf{T}\) , there exists a neighborhood \(\mathbf{V}\) of \(x\) such that
\end{enumerate}

\[
\sup_{\substack{\mathbf{K}\in \mathfrak{K}(\mathbf{T})\\ \mathbf{K}\subset \mathbf{V}}}\operatorname{I}(\mathbf{K}) <   + \infty
\]

(`I is locally bounded').

The measure \(\mu\) is then unique.

The uniqueness of \(\mu\) follows from the Cor. of Prop. 2 of §1, No.~2. The above conditions are necessary, the first three in obvious fashion, the last from the fact that \(\mu^{\bullet}\) is a locally bounded encumbrance, and condition 4) by the Cor. of Prop. 5 of §1, No.~6.

To show that these conditions are sufficient, we begin by treating the case that \(\mathbf{T}\) is compact.

Lemma 1. --- Assume that T is compact, and set \(l = \mathrm{I}(\mathrm{T})\) . For every \(\mathrm{A} \subset \mathrm{T}\) , set

\[
\mathrm{J}(\mathrm{A}) = \sup_{\substack{\mathrm{K}\in \mathfrak{K}(\mathrm{T})\\ \mathrm{K}\subset \mathrm{A}}}\mathrm{I}(\mathrm{K})\tag{1}
\]

and let \(\Phi\) be the set of \(A \subset T\) such that \(J(A) + J(\mathbb{C}A) = l\) . The set \(\Phi\) is then a clan that contains \(\mathfrak{K}(T)\) , and the function \(J\) on \(\Phi\) is increasing and additive.

It is clear that \(J\) is an increasing set function, extending \(I\) , and that \(J(A) + J(\mathbb{C}A) \leqslant l\) for every \(A \subset T\) .

Let K and S be two compact sets in T; we are first going to show that

\[
\mathrm{J} (\mathrm{K} \cap \mathrm{S}) + \mathrm{J} (\mathbb {C} \mathrm{K} \cap \mathrm{S}) = \mathrm{J} (\mathrm{S}).\tag{2}
\]

On considering the restrictions of I to \(\mathfrak{K}(\mathrm{S})\) and of J to \(\mathfrak{P}(\mathrm{S})\) , one reduces immediately to the case that \(\mathrm{S} = \mathrm{T}\) . Since T is normal, K is the intersection of the decreasing directed family of its compact neighborhoods, and condition 4) implies the existence, for every \(\varepsilon > 0\) , of a compact neighborhood H of K such that \(\mathrm{I}(\mathrm{H}) \leqslant \mathrm{I}(\mathrm{K}) + \varepsilon\) . Let L be the closure of T - H; L is compact, \(\mathrm{L} \cap \mathrm{K} = \varnothing\) and \(\mathrm{H} \cup \mathrm{L} = \mathrm{T}\) , therefore \(l = \mathrm{I}(\mathrm{H} \cup \mathrm{L}) \leqslant \mathrm{I}(\mathrm{H}) + \mathrm{I}(\mathrm{L}) \leqslant \mathrm{I}(\mathrm{K}) + \mathrm{I}(\mathrm{L}) + \varepsilon\) (condition 2)), whence the relation \(\mathrm{J}(\mathrm{K}) + \mathrm{J}(\mathbb{C}\mathrm{K}) \geqslant \mathrm{I}(\mathrm{K}) + \mathrm{I}(\mathrm{L}) \geqslant l - \varepsilon\) . Since \(\varepsilon\) is arbitrary, \(\mathrm{J}(\mathrm{K}) + \mathrm{J}(\mathbb{C}\mathrm{K}) = l\) . This proves the formula (2), as well as the inclusion \(\mathfrak{K}(\mathrm{T}) \subset \Phi\) .

Let us now prove that \(\Phi\) is a clan. Since \(\Phi\) is obviously stable under passage to the complement, it suffices to show that if \(A_{1}\) and \(A_{2}\) denote elements of \(\Phi\) , then \(A_{1} \cup A_{2} \in \Phi\) , or again that

\[
\mathrm{J} \left(\mathrm{A} _ {1} \cup \mathrm{A} _ {2}\right) + \mathrm{J} \left(\mathbf {C} \left(\mathrm{A} _ {1} \cup \mathrm{A} _ {2}\right)\right) \geqslant l.\tag{3}
\]

Denote by \(\varepsilon\) a number \(>0\) , and, for \(i = 1,2\) , let \(\mathbf{K}_i\) be a compact set contained in \(\mathbf{A}_i\) , and \(\mathbf{L}_i\) a compact set contained in \(\mathbf{C}\mathbf{A}_i\) , such that

\[
\mathrm{I} \left(\mathrm{K} _ {i}\right) \geqslant \mathrm{J} \left(\mathrm{A} _ {i}\right) - \varepsilon , \quad \mathrm{I} \left(\mathrm{L} _ {i}\right) \geqslant \mathrm{J} \left(\mathbb {C} \mathrm{A} _ {i}\right) - \varepsilon .
\]

Set \(\mathbf{M}_1 = \mathbf{K}_1\cup \mathbf{L}_1\) ; the relations \(l = J(\mathbf{M}_1) + \mathbf{J}(\mathbf{C}\mathbf{M}_1)\) and

\[
\mathrm{J} \left(\mathrm{M} _ {1}\right) = \mathrm{I} \left(\mathrm{K} _ {1}\right) + \mathrm{I} \left(\mathrm{L} _ {1}\right) \geqslant \mathrm{J} \left(\mathrm{A} _ {1}\right) + \mathrm{J} \left(\mathbf {C A} _ {1}\right) - 2 \varepsilon = l - 2 \varepsilon
\]

imply that \(\mathrm{J}(\mathbb{C}\mathrm{M}_1) \leqslant 2\varepsilon\) . Then if \(\mathrm{S}\) is a compact subset of \(\mathrm{T}\) , the relation (2) (applied to \(\mathrm{K} = \mathrm{M}_1\) ) implies \(\mathrm{J}(\mathrm{S}) \leqslant \mathrm{J}(\mathrm{M}_1 \cap \mathrm{S}) + 2\varepsilon\) , whence

\[
\mathrm{J} (\mathrm{S}) \leqslant \mathrm{J} (\mathrm{K} _ {1} \cap \mathrm{S}) + \mathrm{J} (\mathrm{L} _ {1} \cap \mathrm{S}) + 2 \varepsilon .
\]

Let us add the inequalities obtained by making \(S = K_{2}\) and \(S = L_{2}\) , and take into account the inequality \(\mathrm{J}(K_{2}) + \mathrm{J}(L_{2}) \geqslant l - 2\varepsilon\) and the fact that \(K_{1} \cap K_{2}\) , \(L_{1} \cap K_{2}\) and \(K_{1} \cap L_{2}\) are three disjoint compact sets contained in \(A_{1} \cup A_{2}\) . Denoting by C the union of these three compact sets, it follows that

\[
\begin{array}{r l} & {l - 2 \varepsilon \leqslant \mathrm{J} (\mathrm{K} _ {2}) + \mathrm{J} (\mathrm{L} _ {2}) \leqslant \mathrm{J} (\mathrm{C}) + \mathrm{J} (\mathrm{L} _ {1} \cap \mathrm{L} _ {2}) + 4 \varepsilon} \\ & {\qquad \leqslant \mathrm{J} (\mathrm{A} _ {1} \cup \mathrm{A} _ {2}) + \mathrm{J} \big (\mathbb {C} (\mathrm{A} _ {1} \cup \mathrm{A} _ {2}) \big) + 4 \varepsilon ,} \end{array}
\]

whence immediately the desired formula (3), in view of the arbitrariness of \(\varepsilon\) . This having been established, the preceding inequalities imply that \(\mathrm{J}(\mathbf{C}) \geqslant \mathrm{J}(\mathbf{A}_1 \cup \mathbf{A}_2) - 6\varepsilon\) ; if \(\mathbf{A}_1\) and \(\mathbf{A}_2\) are disjoint, then \(\mathbf{C}\) is the union of \(\mathbf{K}_1 \cap \mathbf{L}_2 \subset \mathbf{A}_1\) and \(\mathbf{K}_2 \cap \mathbf{L}_1 \subset \mathbf{A}_2\) , from which one deduces that \(\mathrm{J}(\mathbf{A}_1 \cup \mathbf{A}_2) \leqslant \mathrm{J}(\mathbf{A}_1) + \mathrm{J}(\mathbf{A}_2)\) . The reverse inequality being obvious, \(\mathbf{J}\) is indeed additive on \(\Phi\) , and the lemma is established.

Let us complete the proof of the theorem for the case that T is compact. Let \(\mathcal{E}(\Phi)\) be the vector space of \(\Phi\) -step functions on T, equipped with the topology of uniform convergence (Ch. IV, §4, No.~9, Def. 4); we shall again denote by J the positive linear form on \(\mathcal{E}(\Phi)\) associated with the additive function J (loc. cit., Prop. 18). Since \(J(T) = l\) , J is continuous and has norm \(l\) . Let then \(\mathcal{H}\) be the closure of \(\mathcal{E}(\Phi)\) for the topology of uniform convergence; one verifies at once that J may be extended by continuity to a positive linear form on \(\mathcal{H}\) , again denoted J. Since \(\mathcal{H}\) contains \(\mathcal{C}(T)\) (loc. cit., No.~10, Prop. 19) the restriction of J to \(\mathcal{C}(T)\) is a positive measure \(\mu\) . It remains to show that \(\mu^{\bullet}(K) = I(K)\) for every compact subset K of T. Now, we have \(\mu^{\bullet}(K) = \inf_{f \in S_K} \mu^{\bullet}(f)\) , where \(S_K\) denotes the set of elements of \(\mathcal{C}(T)\) that are \(\geqslant \varphi_K\) (§1, No.~6, Prop. 5). Since \(J(f) = \mu^{\bullet}(f)\) for \(f \in \mathcal{C}(T)\) , it clearly suffices to show that \(J(K) \geqslant \inf_{f \in S_K} J(f)\) . As in the proof of Lemma 1, let H be a compact neighborhood of K such that \(J(H) \leqslant J(K) + \varepsilon\) , and let \(f\) be a continuous function on T, between 0 and 1, equal to 1 on K and to 0 outside H (GT, IX, §4, No.~1, Prop. 1). Then

\[
\mathrm{J} (f) \leqslant \mathrm{J} (\mathrm{H}) \leqslant \mathrm{J} (\mathrm{K}) + \varepsilon ;
\]

\(\varepsilon\) being arbitrary, the desired inequality is proved, and the theorem is thus established when \(\mathbf{T}\) is compact.

Let us now pass to the general case. For every compact set L in T, let \(I_{L}\) be the restriction of I to \(\mathfrak{K}(L)\) . By the special case just treated, there exists a measure \(\mu_{L}\) on L, unique, such that \(\mu_{\mathrm{L}}(\mathrm{K}) = \mathrm{I}_{\mathrm{L}}(\mathrm{K})\) for every compact set \(K \subset L\) . Let then \(L'\) be a compact set contained in L; we have \((\mu_{\mathrm{L}})_{\mathrm{L}'}^{\bullet}(\mathrm{K}) = \mu_{\mathrm{L}}^{\bullet}(\mathrm{K}) = \mu_{\mathrm{L}'}^{\bullet}(\mathrm{K})\) for every compact set \(K \subset L'\) , therefore \(\mu_{\mathrm{L}'} = (\mu_{\mathrm{L}})_{\mathrm{L}'}\) ; the mapping \(\mu : L \mapsto \mu_{L}\) is a premeasure. The condition 5) expresses that \(\mu\) is a measure, and the relation \(\mathrm{I}(\mathrm{K}) = \mu^{\bullet}(\mathrm{K})\) for all compact \(\mathbf{K} \subset \mathbf{T}\) is obvious.

Remarks. --- 1) The condition 4) may be replaced, in the statement of Theorem 1, by the following condition (`right-continuity'):

4') For every \(\mathbf{K} \in \mathfrak{K}(\mathbf{T})\) and every \(\varepsilon > 0\) , there exists an open set \(\mathbf{U}\) containing \(\mathbf{K}\) , such that \(\mathrm{I}(\mathrm{H}) \leqslant \mathrm{I}(\mathrm{K}) + \varepsilon\) for every compact set \(\mathrm{H} \subset \mathbf{U}\) .

For, if \(\mu\) is a measure, the function \(I: K \mapsto \mu^{\bullet}(K)\) satisfies \(4'\) (§1, No.~9, Prop. 13). Conversely, suppose that I satisfies 1) and \(4'\) ; let us show that I then satisfies 4). With notations as in the statement of Theorem 1, choose an \(\varepsilon > 0\) and an open set U containing the compact set \(K = \bigcap_{\alpha \in A} K_{\alpha}\) and such that \(4'\) ) is

satisfied. There then exists an index \(\beta \in A\) such that \(K_{\beta} \subset U\) , and this implies

\[
\inf _ {\alpha \in \mathbf {A}} \mathrm{I} (\mathrm{K} _ {\alpha}) \leqslant \mathrm{I} (\mathrm{K} _ {\beta}) \leqslant \mathrm{I} (\mathrm{K}) + \varepsilon
\]

and 4) is indeed verified.

\begin{enumerate}
\def\labelenumi{\arabic{enumi})}
\setcounter{enumi}{1}
\tightlist
\item
  The set of conditions 2) and 3) may be replaced, in the statement of Theorem 1, by the following condition:
\end{enumerate}

If K and L are compact subsets of T, then

\[
\mathrm{I} (\mathrm{K} \cup \mathrm{L}) + \mathrm{I} (\mathrm{K} \cap \mathrm{L}) = \mathrm{I} (\mathrm{K}) + \mathrm{I} (\mathrm{L}).
\]

Indeed, this condition implies 2) and 3), and on the other hand

\[
\mu^ {\bullet} (\mathrm{K} \cup \mathrm{L}) = \mu^ {\bullet} (\mathrm{K} \cap \mathrm{L}) = \mu^ {\bullet} (\mathrm{K}) + \mu^ {\bullet} (\mathrm{L})
\]

for every measure \(\mu\) , by virtue of the relation \(\varphi_{\mathrm{K}\cup\mathrm{L}} + \varphi_{\mathrm{K}\cap\mathrm{L}} = \varphi_{\mathrm{K}} + \varphi_{\mathrm{L}}\) between the characteristic functions.

\subsection{Inner regular set functions}

DEFINITION 1. --- Let T be a topological space, and let \(\mathfrak{B}(T)\) be the Borel tribe of T; let I be a mapping of \(\mathfrak{B}(T)\) into \(\overline{\mathbf{R}}_{+}\) .

\begin{enumerate}
\def\labelenumi{\alph{enumi})}
\tightlist
\item
  I is said to be countably additive if, for every sequence \((\mathbf{A}_{n})\) of pairwise disjoint elements of \(\mathfrak{B}(\mathbf{T})\) ,
\end{enumerate}

\[
\mathrm{I} \left(\bigcup_ {n} \mathrm{A} _ {n}\right) = \sum_ {n} \mathrm{I} (\mathrm{A} _ {n}).\tag{4}
\]

\begin{enumerate}
\def\labelenumi{\alph{enumi})}
\setcounter{enumi}{1}
\tightlist
\item
  I is said to be inner regular if, for every set \(\mathbf{A} \in \mathfrak{B}(\mathbf{T})\) ,
\end{enumerate}

\[
\mathrm{I} (\mathrm{A}) = \sup _ {\mathrm{K}} \mathrm{I} (\mathrm{K}),\tag{5}
\]

where K runs over the set of compact subsets of A.

\begin{enumerate}
\def\labelenumi{\alph{enumi})}
\setcounter{enumi}{2}
\tightlist
\item
  I is said to be bounded (resp. locally bounded) if \(\mathrm{I(T)} < +\infty\) (resp. if every point \(x \in \mathbf{T}\) admits an open neighborhood V such that \(\mathrm{I(V)} < +\infty\) ).
\end{enumerate}

Remarks. --- 1) The condition a) clearly implies that I is an increasing mapping of \(\mathfrak{B}(\mathbf{T})\) (ordered by inclusion) into \(\overline{\mathbf{R}}_{+}\) .

\begin{enumerate}
\def\labelenumi{\arabic{enumi})}
\setcounter{enumi}{1}
\item
  Suppose that I is countably additive; let \((\mathbf{A}_n)_{n\in \mathbf{N}}\) be an increasing sequence of Borel sets, and let \(\mathbf{A} = \bigcup_{n\in \mathbf{N}}\mathbf{A}_n\) . The sets \(\mathrm{D}_0 = \mathrm{A}_0\) , \(\mathrm{D}_n = \mathrm{A}_n - \mathrm{A}_{n - 1}\) being pairwise disjoint, and their union being A, we have \(\mathrm{I}(\mathrm{A}) = \sum_{n} \mathrm{I}(\mathrm{D}_{n}) = \lim_{n \to \infty} \mathrm{I}(\mathrm{A}_{n})\) . Similarly, if \((\mathrm{B}_{n})\) is a decreasing sequence of Borel sets, and if \(\mathrm{I}(\mathrm{B}_{0}) < +\infty\) , then \(\mathrm{I}\left(\bigcap_{n} \mathrm{B}_{n}\right) = \lim_{n \to \infty} \mathrm{I}(\mathrm{B}_{n})\) : it suffices to apply the preceding to the sets \(A_{n} = B_{0} - B_{n}\) .
\item
  Let \((\mathbf{A}_n)\) be any sequence of Borel sets of T. If I is countably additive, then \(\mathrm{I}\left(\bigcup_{n} \mathrm{A}_n\right) \leqslant \sum_{n} \mathrm{I}(\mathrm{A}_n)\) . By the preceding remark, it suffices to establish this inequality for a finite sequence. One immediately reduces to the case of two sets \(A_{1}\) and \(A_{2}\) ; but the relation (4) implies that
\end{enumerate}

\[
\mathrm{I} \left(\mathrm{A} _ {1} \cup \mathrm{A} _ {2}\right) = \mathrm{I} \left(\mathrm{A} _ {1}\right) + \mathrm{I} \left(\mathrm{A} _ {2} - \left(\mathrm{A} _ {1} \cap \mathrm{A} _ {2}\right)\right) \leqslant \mathrm{I} \left(\mathrm{A} _ {1}\right) + \mathrm{I} \left(\mathrm{A} _ {2}\right).
\]

The desired inequality then follows immediately.

\begin{enumerate}
\def\labelenumi{\arabic{enumi})}
\setcounter{enumi}{3}
\item
  If I is a countably additive and locally bounded function, the preceding remark implies at once that \(\mathrm{I}(\mathbf{K}) < +\infty\) for every compact set \(\mathbf{K} \subset \mathbf{T}\) .
\item
  One can show that if I is additive, that is, satisfies (4) for finite sequences, and if I is inner regular, then I is countably additive (Exer. 7). The reader can also ascertain that only additivity and inner regularity are used in the proof of Th. 2 below.
\end{enumerate}

THEOREM 2. --- Let T be a topological space, and let I be a function defined on \(\mathfrak{B}(T)\) , with values in \(\overline{\mathbf{R}}_{+}\) . In order that there exist a measure \(\mu\) on T such that \(\mu^{\bullet}(\mathbf{A}) = \mathrm{I}(\mathbf{A})\) for every \(\mathbf{A} \in \mathfrak{B}(T)\) , it is necessary and sufficient that I be countably additive, locally bounded and inner regular. The measure \(\mu\) is then unique.

These three conditions are necessary: for, the mapping \(A \mapsto \mu^{\bullet}(A)\) on \(\mathfrak{B}(T)\) is countably additive ( \(\S1\) , No.~5, Cor. of Prop. 4), locally bounded by the definition of measures ( \(\S1\) , No.~2, Def. 5), and inner regular by Remark 3 of \(\S1\) , No.~2.

We pass to existence. It is clear that the restriction of I to \(\mathcal{K}(\mathrm{T})\) satisfies conditions 1), 2), 3) and 5) of the statement of Th. 1; let us show that 4) is satisfied as well. Let K be a compact subset of T, the intersection of a decreasing directed family \((\mathrm{K}_{\alpha})_{\alpha \in \mathrm{A}}\) of compact sets, and let \(\varepsilon\) be a number \(>0\) ; I being locally bounded, there exists an open (hence Borel) neighborhood V of K such that \(\mathrm{I(V)} < +\infty\) , and then there exists an index \(\alpha\) such that \(\mathrm{K}_{\alpha} \subset \mathrm{V}\) ; changing notation if necessary, we can suppose that \(\mathrm{K}_{\alpha} \subset \mathrm{V}\) for all \(\alpha \in \mathrm{A}\) . By the inner regularity of I, there exists a compact set \(\mathrm{L} \subset \mathrm{V} - \mathrm{K}\) such that \(\mathrm{I(L)} \geqslant \mathrm{I(V} - \mathrm{K)} - \varepsilon\) ; since L does not intersect K, there exists an index \(\alpha\) such that \(\mathrm{L} \cap \mathrm{K}_{\alpha} = \varnothing\) , and one then has \(\mathrm{I}(\mathrm{V} - \mathrm{K}_{\alpha}) \geqslant \mathrm{I}(\mathrm{L}) \geqslant \mathrm{I}(\mathrm{V} - \mathrm{K}) - \varepsilon\) . Since \(\mathbf{K}_{\alpha} \subset \mathbf{V}\) , it follows that \(\mathrm{I}(\mathbf{K}_{\alpha}) \leqslant \mathrm{I}(\mathbf{K}) + \varepsilon\) and the condition 4) is verified.

By Th. 1, there exists a measure \(\mu\) such that \(\mu^{\bullet}(\mathrm{K}) = \mathrm{I}(\mathrm{K})\) for all \(\mathbf{K} \in \mathfrak{K}(\mathbf{T})\) . The inner regularity of the set functions \(\mu^{\bullet}\) and I on \(\mathfrak{B}(\mathbf{T})\) then implies that \(\mu^{\bullet}(\mathbf{A}) = \mathrm{I}(\mathbf{A})\) for all \(\mathbf{A} \in \mathfrak{B}(\mathbf{T})\) , and existence is proved. The uniqueness of \(\mu\) follows from the uniqueness assertion of Th. 1.

\subsection{Radon spaces}

DEFINITION 2. --- Let T be a topological space. One says that T is a Radon (resp. strongly Radon) space if T is Hausdorff and if every function defined on the Borel tribe \(\mathfrak{B}(\mathrm{T})\) of T, with values in \(\overline{R}_{+}\) , that is countably additive and bounded (resp. locally bounded) is inner regular.

For example, we shall see later on (Prop. 3) that every Polish space is strongly Radon. In particular, every locally compact space with a countable base is strongly Radon.

There exist Radon spaces that are not strongly Radon.

PROPOSITION 1. --- Every Lindelöf \(^{(1)}\) Radon space is strongly Radon.

Let T be a Lindelöf space that is Radon, and let I be a set function on the tribe \(\mathfrak{B}(T)\) that is positive, countably additive and locally bounded. The open sets V such that \(\mathrm{I}(\mathrm{V}) < +\infty\) form a covering of T, from which one can extract a countable covering \((\mathrm{V}_{n})_{n \in \mathbb{N}}\) . Set \(G_{n} = V_{0} \cup V_{1} \cup \cdots \cup V_{n}\) for every \(n \in N\) ; set \(H_{0} = G_{0}\) and \(H_{n} = G_{n} - G_{n-1}\) for \(n \geqslant 1\) ; finally, denote by \(I_{n}\) the set function \(A \mapsto \mathrm{I}(A \cap H_{n})\) on \(\mathfrak{B}(T)\) , which is obviously countably additive and bounded. Since the sets \(H_{n}\) form a partition of T, we have \(I = \sum_{n} I_{n}\) . The space T being Radon, for every \(n \in N\) there exists a bounded measure \(\mu_{n}\) on T such that \(\mu_{n}^{\bullet}(\mathrm{A}) = I_{n}(\mathrm{A})\) for all \(A \in \mathfrak{B}(T)\) ; therefore also \(\sum_{n} \mu_{n}^{\bullet}(\mathrm{A}) = I(\mathrm{A})\) . Since I is locally bounded, the family \((\mu_{n})\) is summable (§1, No.~7, Prop. 7); if \(\mu\) denotes \(\sum_{n} \mu_{n}\) , we have \(\mu^{\bullet}(\mathrm{A}) = I(\mathrm{A})\) for all \(A \in \mathfrak{B}(T)\) , and it follows that I is inner regular. In other words, T is strongly Radon.

Recall that a subset A of a topological space T is said to be universally measurable if A is \(\mu\) -measurable for every measure \(\mu\) on T. This is equivalent to saying that A is \(\mu\) -measurable for every measure \(\mu\) on T with compact support ( \(\S1\) , No.~8, Prop. 9).

PROPOSITION 2. --- Let X be a topological space and T a subspace of X.

\begin{enumerate}
\def\labelenumi{\alph{enumi})}
\tightlist
\item
  Suppose that \(\mathbf{T}\) is a Radon space. Then, for every function I defined on \(\mathfrak{B}(\mathbf{X})\) that is positive, countably additive and bounded, one has
\end{enumerate}

\[
\sup_{\substack{K\text{compact}\\ K\subset T}}I(K) = \inf_{\substack{B\in \mathfrak{B}(X)\\ B\supset T}}I(B).\tag{6}
\]

Moreover, \(\mathbf{T}\) is universally measurable in \(\mathbf{X}\) .

\begin{enumerate}
\def\labelenumi{\alph{enumi})}
\setcounter{enumi}{1}
\tightlist
\item
  Conversely, suppose that X is a Radon space and that T is universally measurable in X; then T is a Radon space.
\end{enumerate}

Let us prove \(a\) ). We denote by \(\alpha\) the right side of (6); for every \(n \in \mathbf{N}\) , there exists a set \(C_n \in \mathfrak{B}(X)\) containing \(T\) , such that \(I(C_n) \leqslant \alpha + 2^{-n}\) . Setting \(C = \bigcap_{n} C_n\) , we then have \(T \subset C\) , \(I(C) = \alpha\) . If \(A \in \mathfrak{B}(T)\) , let us choose a Borel set \(B\) of \(X\) such that \(A = B \cap T\) (GT, IX, §6, No.~3) and set \(J(A) = I(B \cap C)\) . This number does not depend on the choice of \(B\) , for if \(B'\) is a second Borel set in \(X\) such that \(A = B' \cap T\) , then \(B \cap C\) and \(B' \cap C\) differ only by a Borel set \(M\) contained in \(C - T\) , and \(I(M) = 0\) by the construction of \(C\) . Clearly \(J(K) = I(K)\) for every compact set \(K \subset T\) . Let \((A_n)\) be a sequence of Borel sets of \(T\) , pairwise disjoint, and, for each \(n\) , let \(B_n\) be a Borel set of \(X\) such that \(B_n \cap T = A_n\) . Replacing \(B_n\) by \(B_n - \left( \bigcup_{k < n} B_k \right)\) if necessary, we can suppose that the sets \(B_n\) are pairwise disjoint. Set \(A = \bigcup_{n} A_n\) and \(B = \bigcup_{n} B_n\) ; then

\[
\mathrm{J} (\mathrm{A}) = \mathrm{I} (\mathrm{B} \cap \mathrm{C}) = \sum_ {n} \mathrm{I} (\mathrm{B} _ {n} \cap \mathrm{C}) = \sum_ {n} \mathrm{J} (\mathrm{A} _ {n});
\]

J is thus a countably additive and bounded function on \(\mathfrak{B}(\mathrm{T})\) . Since T is by hypothesis a Radon space, there exists a bounded measure \(\mu\) on T such that \(\mathrm{J}(\mathrm{A}) = \mu^{\bullet}(\mathrm{A})\) for every \(\mathrm{A} \in \mathfrak{B}(\mathrm{T})\) ; consequently

\[
\alpha = \mathrm{J} (\mathrm{T}) = \mu^ {\bullet} (\mathrm{T}) = \sup _ {\mathrm{K}} \mu^ {\bullet} (\mathrm{K}) = \sup _ {\mathrm{K}} \mathrm{J} (\mathrm{K}),
\]

by the definition of \(\mu^{\bullet}\) . Formula (6) is thus established.

Let us now show that T is universally measurable. Let \(\lambda\) be a bounded measure on X; the preceding argument may be applied to the set function \(I: A \mapsto \lambda^{\bullet}(A)\) on \(\mathfrak{B}(X)\) , thus there exists a sequence \((K_{n})\) of compact subsets of T such that (with the above notations)

\[
\sup _ {n} \lambda^ {\bullet} (\mathrm{K} _ {n}) = \mathrm{J} (\mathrm{T}) = \lambda^ {\bullet} (\mathrm{C}).
\]

Set \(K' = \bigcup_{n \in \mathbf{N}} K_n\) ; \(K'\) is Borel in \(X\) , \(K' \subset T \subset C\) , \(\lambda^\bullet(K') = \lambda^\bullet(C)\) , therefore these three sets differ only by \(\lambda\) -negligible sets, and so \(T\) is \(\lambda\) -measurable. This completes the proof of \(a\) ).

Let us pass to \(b\) ). Suppose that \(X\) is a Radon space, and that \(T\) is universally measurable in \(X\) . Let \(I\) be a positive function on \(\mathfrak{B}(T)\) that is countably additive and bounded; the function \(A \mapsto I(A \cap T)\) on \(\mathfrak{B}(X)\) is then positive, countably additive and bounded, therefore there exists a bounded measure \(\nu\) on \(X\) such that \(I(A \cap T) = \nu^{\bullet}(A)\) for all \(A \in \mathfrak{B}(X)\) . Now, \(T\) is \(\nu\) -measurable; the preceding relation shows that \(\nu^{\bullet}(K) = 0\) for every compact subset \(K\) of \(X\) that is disjoint from \(T\) , therefore \(\nu\) is concentrated on \(T\) . Consequently, for every Borel set \(A\) of \(X\) , we have \(I(A \cap T) = \nu^{\bullet}(A \cap T) = \mu^{\bullet}(A \cap T)\) , where \(\mu\) is the measure induced by \(\nu\) on \(T\) . Finally, it follows that \(I(B) = \mu^{\bullet}(B)\) for every set \(B \in \mathfrak{B}(T)\) (GT, IX, §6, No.~3, Remark 2), and \(I\) is indeed inner regular.

COROLLARY. --- If X is a Radon space, then every Borel subset T of X is Radon.

For, \(\mathbf{T}\) is universally measurable in \(\mathbf{X}\) .

PROPOSITION 3. --- Every Souslin space (in particular, every Polish or Lusin space) is strongly Radon.

Let T be a Souslin space; since T is a Lindelöf space (TG, IX, Appendix I, Cor. of Prop. 1), \(^{(2)}\) it suffices to show that T is Radon (Prop. 1). Let I be a function defined on \(\mathfrak{B}(T)\) , positive, countably additive and bounded. We extend I to \(\mathfrak{P}(T)\) by setting, for every subset A of T,

\[
\operatorname{I}(\mathbf{A}) = \inf_{\substack{\mathbf{B}\in \mathfrak{B}(\mathbf{T})\\ \mathbf{B}\supset \mathbf{A}}}\operatorname{I}(\mathbf{B}).
\]

Let us show that this extension is a capacity on T (GT, IX, §6, No.~9). It is clear that the relation \(\mathbf{A} \subset \mathbf{A}'\) implies \(\mathrm{I}(\mathbf{A}) \leqslant \mathrm{I}(\mathbf{A}')\) . Let \((\mathbf{A}_n)\) be an increasing sequence of subsets of T, and let \(\mathbf{A} = \bigcup_{n} \mathbf{A}_n\) . The set of Borel sets that contain \(\mathbf{A}_n\) being stable for countable intersections, there exists for each \(n\) a Borel set \(\mathbf{B}_n\) such that \(\mathbf{A}_n \subset \mathbf{B}_n\) and \(\mathrm{I}(\mathbf{A}_n) = \mathrm{I}(\mathbf{B}_n)\) (cf.~the proof of Prop. 2). Set \(\mathbf{C}_n = \bigcap_{p \geqslant n} \mathbf{B}_p\) ; \(\mathbf{C}_n\) is Borel, and \(\mathbf{A}_n \subset \mathbf{C}_n \subset \mathbf{B}_n\) , therefore \(\mathrm{I}(\mathbf{A}_n) = \mathrm{I}(\mathbf{C}_n)\) . On the other hand, the sequence \((\mathbf{C}_n)\) is increasing. Let \(\mathbf{C} = \bigcup_{n} \mathbf{C}_n\) : the relation \(\mathbf{A} \subset \mathbf{C}\) implies that

\[
\mathrm{I} (\mathrm{A}) \leqslant \mathrm{I} (\mathrm{C}) = \lim _ {n} \mathrm{I} (\mathrm{C} _ {n}) = \lim _ {n} \mathrm{I} (\mathrm{A} _ {n}),
\]

whence the equality \(\mathrm{I}(\mathrm{A})=\lim_{n}\mathrm{I}(\mathrm{A}_{n})\) is immediate. Consequently, I is a capacity.

If \((\mathbf{H}_{n})\) is a decreasing sequence of closed sets in T, obviously \(\mathrm{I}\left(\bigcap_{n}\mathrm{H}_{n}\right)=\inf_{n}\mathrm{I}(\mathrm{H}_{n})\) . It follows that every Souslin subset F of T is capacitable for I (TG, IX, §6, No.~10, Prop. 15). In particular, every Borel set A of T is capacitable (loc. cit., §6, No.~3, Prop. 10). \(^{(3)}\) In other words,

\[
\mathrm{I} (\mathrm{A}) = \sup _ {\mathrm{K}} \mathrm{I} (\mathrm{K}),
\]

where K runs over the set of compact sets contained in A; we have proved that I is inner regular.

Remark. --- Let X be a Lusin space (in particular, any Polish space), and f a bijective continuous mapping of X onto a (Lusin) regular space Y. One knows (TG, IX, §6, No.~7, Prop. 14) that the mapping \(B \mapsto f^{-1}(B)\) is a bijection of the Borel tribe of Y onto the Borel tribe of X. The spaces X and Y are Lusin, hence strongly Radon (Prop. 3). It follows immediately that the mapping \(\mu \mapsto f(\mu)\) is a bijection of the set of bounded measures on X onto the set of bounded measures on Y. \(^{(4)}\)

\section{INVERSE LIMITS OF MEASURES}

Throughout this section, I denotes a nonempty set, equipped with a preorder relation, denoted \(i \leqslant j\) , and directed for this relation. Recall (GT, I, §4, No.~4) that an inverse system of topological spaces indexed by I is a family \((\mathrm{T}_i, p_{ij})\) where \(\mathrm{T}_i\) is a topological space and \(p_{ij}\) is a continuous mapping of \(\mathrm{T}_j\) into \(\mathrm{T}_i\) for \(i \leqslant j\) , where \(p_{ii}\) is the identity mapping of \(\mathrm{T}_i\) , and where \(p_{ik} = p_{ij} \circ p_{jk}\) for \(i \leqslant j \leqslant k\) . Let T be a topological space and \((p_i)_{i \in I}\) a family of continuous mappings \(p_i: \mathrm{T} \to \mathrm{T}_i\) . The family \((p_i)_{i \in I}\) is said to be coherent if \(p_i = p_{ij} \circ p_j\) for \(i \leqslant j\) , and it is said to be separating if for distinct \(x, y\) in T there exists an \(i \in I\) such that \(p_i(x) \neq p_i(y)\) . When \(\mathrm{T} = \varprojlim \mathrm{T}_i\) and \(p_i\) is the canonical mapping of T into \(\mathrm{T}_i\) , the family \((p_i)_{i \in I}\) is coherent and separating.

\subsection{Complements on compact spaces and inverse limits}

PROPOSITION 1. --- Let X and Y be two topological spaces and f a continuous mapping of X into Y. Let \((\mathrm{K}_{\alpha})_{\alpha \in \mathrm{A}}\) be a decreasing directed family of compact subsets of X, with intersection K. Then \(f(\mathrm{K}) = \bigcap f(\mathrm{K}_{\alpha})\) .

For, let \(y\) be a point of \(\bigcap_{\alpha \in \mathbf{A}} f(\mathbf{K}_{\alpha})\) ; for every \(\alpha \in \mathbf{A}\) , the set \(\mathbf{L}_{\alpha} = \mathbf{K}_{\alpha} \cap \overline{f}^{-1}(y)\) is compact and nonempty. The family \((\mathbf{L}_{\alpha})_{\alpha \in \mathbf{A}}\) is directed downward, therefore its intersection \(\mathbf{L}\) is nonempty. Now, \(\mathbf{L} = \mathbf{K} \cap \overline{f}^{-1}(y)\) , whence \(y \in f(\mathbf{K})\) . We have thus proved the inclusion \(f(\mathbf{K}) \supset \bigcap_{\alpha \in \mathbf{A}} f(\mathbf{K}_{\alpha})\) , and the reverse inclusion is obvious.

PROPOSITION 2. --- Let there be given an inverse system \((\mathrm{T}_{i}, p_{ij})\) of topological spaces indexed by I, a topological space T, and a coherent and separating family of continuous mappings \(p_{i}: T \to T_{i}\) . Then:

\begin{enumerate}
\def\labelenumi{\alph{enumi})}
\item
  For every compact subset K of T, one has \(\mathbf{K} = \bigcap_{i\in I}\overline{p_i}^{-1}\big(p_i(\mathbf{K})\big)\) .
\item
  Let K and L be two disjoint compact subsets of T. There exists an \(i \in I\) such that \(p_{j}(K)\) and \(p_{j}(L)\) are disjoint for \(j \geqslant i\) .
\item
  Let \(x\) be a point of \(\bigcap_{i \in I} \overline{p_i}^{-1}(p_i(K))\) ; for every \(i \in I\) , the set \(K_i\) of points \(y\) in \(K\) such that \(p_i(y) = p_i(x)\) is a nonempty closed subset of \(K\) . For \(i \leqslant j\) we have \(K_i \supset K_j\) , and, since \(K\) is compact, the set \(\bigcap_{i \in I} K_i\) is therefore nonempty. Let \(y\) be a point of \(\bigcap_{i \in I} K_i\) ; we have \(y \in K\) and \(p_i(y) = p_i(x)\) for all \(i \in I\) , whence \(y = x\) ; finally, \(x \in K\) , which proves the inclusion \(K \supset \bigcap_{i \in I} \overline{p_i}^{-1}(p_i(K))\) ; the reverse inclusion is obvious.
\item
  For every \(i \in \mathbf{I}\) , set \(\mathbf{M}_i = \overline{p_i}^{-1}(p_i(\mathbf{K})) \cap \mathbf{L}\) ; this is a closed subset of the compact space \(\mathbf{L}\) , we have \(\mathbf{M}_i \supset \mathbf{M}_j\) for \(i \leqslant j\) , and, by \(a\) ),
\end{enumerate}

\[
\bigcap_ {i \in I} M _ {i} = K \cap L = \varnothing .
\]

Consequently, there exists an index \(i\) such that \(\mathbf{M}_i = \varnothing\) . For \(j \geqslant i\) , we have \(\overline{p_j}^{-1} (p_j(\mathbf{K})) \cap \mathbf{L} = \mathbf{M}_j \subset \mathbf{M}_i = \varnothing\) , whence \(p_j(\mathbf{K}) \cap p_j(\mathbf{L}) = \varnothing\) .
